% ============================================================
% chapter5.tex -- Chapter 5 wrapper.
%
% chapter5_C5D.tex is a SECTION (\section{...}), not a chapter. It needs a
% \chapter{} above it or LaTeX has nowhere to attach it.
%
% Add to thesis_main.tex, in the document body, after chapter 4:
%
%     % ============================================================
% chapter5.tex -- Chapter 5 wrapper.
%
% chapter5_C5D.tex is a SECTION (\section{...}), not a chapter. It needs a
% \chapter{} above it or LaTeX has nowhere to attach it.
%
% Add to thesis_main.tex, in the document body, after chapter 4:
%
%     % ============================================================
% chapter5.tex -- Chapter 5 wrapper.
%
% chapter5_C5D.tex is a SECTION (\section{...}), not a chapter. It needs a
% \chapter{} above it or LaTeX has nowhere to attach it.
%
% Add to thesis_main.tex, in the document body, after chapter 4:
%
%     % ============================================================
% chapter5.tex -- Chapter 5 wrapper.
%
% chapter5_C5D.tex is a SECTION (\section{...}), not a chapter. It needs a
% \chapter{} above it or LaTeX has nowhere to attach it.
%
% Add to thesis_main.tex, in the document body, after chapter 4:
%
%     \include{chapter5}          % if the other chapters use \include
%   or
%     \input{chapter5}            % if they use \input
%
% Use whichever the neighbouring chapters use -- mixing them causes the page
% breaks and numbering to differ from the rest of the thesis.
% ============================================================

\chapter{What the Model Says}
\label{ch:results}

% ------------------------------------------------------------------
% PLACEHOLDER so the chapter is never empty while sections are drafted.
% Delete this block once sec:error_maps exists.
% ------------------------------------------------------------------
\section{The error maps}
\label{sec:error_maps}

\emph{[TO BE WRITTEN --- C5-F. Defines the operating grid, presents
$\epsstep$ and $\epsplat$ over it, and gives the headline numbers with their
scopes. Section~\ref{sec:mechanism} below is drafted and refers back to this
section for the maps it explains.]}

% ------------------------------------------------------------------
% C5-D, drafted 24 August 2026.
% ------------------------------------------------------------------
\input{chapter5_C5D}

% ------------------------------------------------------------------
% Sections still to be written. The labels are declared here so that the
% cross-references in chapter5_C5D.tex resolve rather than printing "??".
% Replace each stub with the real section as it is drafted.
% ------------------------------------------------------------------
\section{Does the transient visit the plateau?}
\label{sec:bridge}

\emph{[TO BE WRITTEN --- C5-A writeup. Full 43-state integration; the
dynamic bridge; the closed form for the early-response slope $k$, which was
cut from Section~\ref{sec:mechanism} because it needs the transient
established first.]}

\section{How long the error persists}
\label{sec:persistence}

\emph{[TO BE WRITTEN. $\tau_{\rm eff}/\tauqss = 1.01$ and $1.07$.]}

\section{Optical thickness and the restricted range}
\label{sec:optical_depth}

\emph{[TO BE WRITTEN --- may instead live in Chapter~\ref{ch:theory};
Section~\ref{sec:mechanism} refers to it for the $\Te \ge 2$~eV restriction.]}
          % if the other chapters use \include
%   or
%     % ============================================================
% chapter5.tex -- Chapter 5 wrapper.
%
% chapter5_C5D.tex is a SECTION (\section{...}), not a chapter. It needs a
% \chapter{} above it or LaTeX has nowhere to attach it.
%
% Add to thesis_main.tex, in the document body, after chapter 4:
%
%     \include{chapter5}          % if the other chapters use \include
%   or
%     \input{chapter5}            % if they use \input
%
% Use whichever the neighbouring chapters use -- mixing them causes the page
% breaks and numbering to differ from the rest of the thesis.
% ============================================================

\chapter{What the Model Says}
\label{ch:results}

% ------------------------------------------------------------------
% PLACEHOLDER so the chapter is never empty while sections are drafted.
% Delete this block once sec:error_maps exists.
% ------------------------------------------------------------------
\section{The error maps}
\label{sec:error_maps}

\emph{[TO BE WRITTEN --- C5-F. Defines the operating grid, presents
$\epsstep$ and $\epsplat$ over it, and gives the headline numbers with their
scopes. Section~\ref{sec:mechanism} below is drafted and refers back to this
section for the maps it explains.]}

% ------------------------------------------------------------------
% C5-D, drafted 24 August 2026.
% ------------------------------------------------------------------
\input{chapter5_C5D}

% ------------------------------------------------------------------
% Sections still to be written. The labels are declared here so that the
% cross-references in chapter5_C5D.tex resolve rather than printing "??".
% Replace each stub with the real section as it is drafted.
% ------------------------------------------------------------------
\section{Does the transient visit the plateau?}
\label{sec:bridge}

\emph{[TO BE WRITTEN --- C5-A writeup. Full 43-state integration; the
dynamic bridge; the closed form for the early-response slope $k$, which was
cut from Section~\ref{sec:mechanism} because it needs the transient
established first.]}

\section{How long the error persists}
\label{sec:persistence}

\emph{[TO BE WRITTEN. $\tau_{\rm eff}/\tauqss = 1.01$ and $1.07$.]}

\section{Optical thickness and the restricted range}
\label{sec:optical_depth}

\emph{[TO BE WRITTEN --- may instead live in Chapter~\ref{ch:theory};
Section~\ref{sec:mechanism} refers to it for the $\Te \ge 2$~eV restriction.]}
            % if they use \input
%
% Use whichever the neighbouring chapters use -- mixing them causes the page
% breaks and numbering to differ from the rest of the thesis.
% ============================================================

\chapter{What the Model Says}
\label{ch:results}

% ------------------------------------------------------------------
% PLACEHOLDER so the chapter is never empty while sections are drafted.
% Delete this block once sec:error_maps exists.
% ------------------------------------------------------------------
\section{The error maps}
\label{sec:error_maps}

\emph{[TO BE WRITTEN --- C5-F. Defines the operating grid, presents
$\epsstep$ and $\epsplat$ over it, and gives the headline numbers with their
scopes. Section~\ref{sec:mechanism} below is drafted and refers back to this
section for the maps it explains.]}

% ------------------------------------------------------------------
% C5-D, drafted 24 August 2026.
% ------------------------------------------------------------------
% ============================================================
% Chapter 5, Section: the mechanism.  C5-D, draft 2.
%
% Written under the thesis-writing skill, 24 August 2026.
% Source: CH5_EVIDENCE.md. L_grid SHA 2d92b58e..., S_grid SHA 7822f536...
% Draft 1 retained as chapter5_C5D_draft1_20260824.tex (failed six gates).
%
% MACROS beyond those chapter3.tex exercises:  \epsstep  \epsplat
%   (recorded in notation_and_definitions.md Sec. 7, never yet compiled --
%    CONFIRM they exist in thesis_main.tex before building.)
% ============================================================

%-----------------------------------------------------------------------------
\section{Why the error is largest where it is}
\label{sec:mechanism}
%-----------------------------------------------------------------------------

The maps of Section~\ref{sec:error_maps} show where the diagnostic error is
large. They do not say why, and a map without a mechanism is a lookup table
that cannot be extrapolated past the grid that produced it. This section asks
which physical quantity puts the ridge where it is.

The answer has two halves and they disagree. The structure along the density
axis and the structure along the temperature axis have different causes, and
reading them as one effect --- which the appearance of a single ridged surface
invites --- gets both wrong.

\subsection{Two error measures}
\label{sec:eps_definitions}

Section~\ref{sec:two_references} fixed the two reference states against which an
observable can be judged. Two further quantities are needed here, and both are
properties of a temperature \emph{step} rather than of a single operator, which
is why neither appears in Chapter~\ref{ch:theory}.

Let the plasma sit at the collisional--radiative equilibrium of
$\Lmat(\Te,\nel)$, and let the electron temperature step to $\Te+\Delta\Te$ at
$t=0$ with $\nel$ held fixed. Write $\Lmat^-$ and $\Lmat^+$ for the operators
before and after, in the notation of Section~\ref{sec:two_clocks}. The
\emph{step error}
\begin{equation}
  \epsstep = \left|\frac{\Rcre(\Te,\nel)}
                        {\Rcre(\Te+\Delta\Te,\nel)} - 1\right|
  \label{eq:eps_step_def}
\end{equation}

\noindent is the dimensionless distance between the two tabulated answers: what
the reading would change by if the plasma re-equilibrated instantly. The
\emph{plateau error}
\begin{equation}
  \epsplat = \left|\frac{\Robs^{\rm PE}}
                        {\Rcre(\Te+\Delta\Te,\nel)} - 1\right|,
  \qquad
  \Robs^{\rm PE} \equiv \Rqss\!\left(\Te+\Delta\Te,\nel;\,\bdep^-\right),
  \label{eq:eps_plateau_def}
\end{equation}

\noindent is the distance between the tabulated answer and the state the plasma
occupies once the excited manifold has settled but before the ground state has
moved. The superscript PE denotes \emph{partial equilibrium}: the post-step
excited response of Eq.~\eqref{eq:qss_ratio}, evaluated on the \emph{pre-step}
reservoir.

The consequence of keeping these separate is that they are not close. Across the
operating grid $\epsplat$ exceeds $\epsstep$ at every point tested, typically by
an order of magnitude (Section~\ref{sec:error_maps}). A diagnostician who
pictures the reading moving smoothly between the two tabulated values has not
merely misjudged the size of the excursion but has the wrong shape of transient.
$\epsplat$ is the quantity this chapter reports, and Section~\ref{sec:bridge}
establishes that the time-dependent system does pass through the state it
measures.

\subsection{An exact decomposition of the finite step}
\label{sec:exact_decomposition}

Chapter~\ref{ch:theory} worked in the ground-state departure coefficient
$\bdep = \ngs/(\Zsb\nel\nion)$, because that is the variable in which the
observable's response has the closed form of Eq.~\eqref{eq:R_of_b}. Here it is
better to use the reservoir itself,
\begin{equation}
  u \;\equiv\; \frac{\ngs}{\nion} \;=\; \bdep\,\Zsb(\Te)\,\nel ,
  \label{eq:u_def}
\end{equation}

\noindent a dimensionless ground-state population per unit ion density. At
fixed $(\Te,\nel)$ the two are proportional, so $\mathrm{d}\ln u =
\mathrm{d}\ln\bdep$ and the sensitivity of
Section~\ref{sec:ground_fed_fraction} carries over unchanged. The reason for
the change of variable is that a temperature step moves $\Zsb$: comparing a
pre-step $\bdep$ with a post-step one would mix a change in the reservoir with
a change in the Saha--Boltzmann normalisation, and $u$ has no normalisation to
move. That confusion is not hypothetical --- an earlier version of the analysis
reported in this section made exactly it, and the resulting contamination in
$\Delta\ln\bdep$ was comparable to the quantity being measured.

\textbf{Assumption.} Equation~\eqref{eq:u_def} treats $\nion$ as a fixed
reservoir, unchanged by the step, as throughout this thesis
(Section~\ref{sec:rate_equation}). This removes the freedom for the ionisation
balance to respond, and it is what makes $u$ numerically equal to the
ground-state population in the solved system. Chapter~\ref{ch:limits} returns
to what a mobile ion population would change.

With that, Eq.~\eqref{eq:R_of_b} makes the post-step observable a ratio of two
affine functions of $u$, so the logarithmic sensitivity of
Section~\ref{sec:ground_fed_fraction} is $\mathrm{d}\ln\Robs/\mathrm{d}\ln u =
\ffrac{3}-\ffrac{4}$, the difference of the two ground-fed fractions.
Integrating between the frozen and the equilibrated reservoir gives, with no
approximation,
\begin{equation}
  \ln\frac{\Robs^{\rm PE}}{\Rcre}
  \;=\; \int_{\ln u^+}^{\ln u^-}\!\!\big(\ffrac{3}-\ffrac{4}\big)\,\mathrm{d}\ln u
  \;=\; \overline{S}\;\Delta\ln u ,
  \label{eq:exact_decomposition}
\end{equation}

\noindent where $\Delta\ln u = \ln(u^-/u^+)$ is the reservoir displacement and
\begin{equation}
  \overline{S} \;\equiv\;
  \frac{1}{\Delta\ln u}\int_{\ln u^+}^{\ln u^-}\!\!
  \big(\ffrac{3}-\ffrac{4}\big)\,\mathrm{d}\ln u
  \label{eq:Sbar_def}
\end{equation}

\noindent is the mean sensitivity over the interval the reservoir traverses.
Both factors are dimensionless. Since $\ffrac{3},\ffrac{4}\in[0,1]$ by
Section~\ref{sec:ground_fed_fraction}, and $\overline{S}$ is an average of their
difference, $|\overline{S}|$ is bounded by the peak sensitivity of
Eq.~\eqref{eq:peak_sensitivity} and so is strictly less than unity everywhere.
The plateau error follows:
\begin{equation}
  \boxed{\ \epsplat = \left|\,\mathrm{e}^{\overline{S}\,\Delta\ln u} - 1\,\right|\ }
  \label{eq:eps_plateau_exact}
\end{equation}

The consequence worth extracting is that a single scalar error has become a
product of two quantities that can be measured independently: how strongly the
observable responds to its reservoir, and how far the reservoir moved. Nothing
in Eq.~\eqref{eq:exact_decomposition} says which of them varies across the
operating map, and that is what the rest of this section determines.

\textbf{Use Eq.~\eqref{eq:eps_plateau_exact}, not its linearisation.} The
first-order form $\epsplat \approx |\overline{S}\,\Delta\ln u|$ is adequate for
a small step and wrong for a large one: at the four-interval step used below it
returns $59.4\,\%$ where the exact expression gives $81.1\,\%$, understating by
more than a quarter (\texttt{verify\_ridge\_mechanism.py}). The exponential
costs nothing.

One caution about what follows. Because Eq.~\eqref{eq:exact_decomposition} is an
identity, comparing its two sides tests nothing: a check that $\epsplat$ equals
$|\mathrm{e}^{\overline{S}\Delta\ln u}-1|$ can fail only through arithmetic
error. What the identity permits is an \emph{attribution} --- the two factors
are separately measurable, so one may ask which carries the structure the maps
show. That is a weaker kind of statement than a derivation, and it is the kind
made below.

\subsection{Across density: a sensitivity structure}
\label{sec:density_mechanism}

The plateau error has an interior maximum in density at
$\nel \approx 1.93\times10^{13}\,\si{\per\cubic\centi\metre}$, at every
temperature in the range this thesis defends.

Evaluating both factors of Eq.~\eqref{eq:exact_decomposition} at every grid
point settles which of them puts it there
(\texttt{verify\_ridge\_mechanism.py}, over $\Te = 2$--$9\,$eV and all eight
density columns). The column maximising $\epsplat$ is the column maximising
$|\overline{S}|$ in $32$ of $32$ temperature rows for a one-interval
temperature step and $31$ of $31$ for a four-interval step; the column
maximising $|\Delta\ln u|$ coincides with the error maximum in $0$ of $32$ and
$0$ of $31$.

Coincidence counts are weak evidence on their own, because a quantity that
barely varies still has its maximum somewhere. The amplitudes settle it. Across
the three decades of density $|\overline{S}|$ varies by a factor $7.11$ and
$|\Delta\ln u|$ by a factor $1.01$. \textbf{A quantity that changes by one per
cent cannot produce a ridge in which the error changes by a factor of seven},
wherever its own maximum falls. Table~\ref{tab:mechanism_profiles} makes the
point without statistics.

\begin{table}[htbp]
\centering
\caption{Density profiles of the plateau error and of the two factors of
Eq.~\eqref{eq:exact_decomposition}, each normalised to its own row maximum and
then taken as the median over the $32$ temperature rows with
$\Te = 2$--$9\,$eV, for a one-interval temperature step. From
\texttt{verify\_ridge\_mechanism.py}.}
\label{tab:mechanism_profiles}
\sisetup{table-format=1.3}
\begin{tabular}{lSSSSSSSS}
  \toprule
  $\nel$ (\si{\per\cubic\centi\metre}) & {$10^{12}$} & {$2.7{\cdot}10^{12}$} &
  {$7.2{\cdot}10^{12}$} & {$1.9{\cdot}10^{13}$} & {$5.2{\cdot}10^{13}$} &
  {$1.4{\cdot}10^{14}$} & {$3.7{\cdot}10^{14}$} & {$10^{15}$} \\
  \midrule
  $\epsplat$        & 0.437 & 0.683 & 0.937 & 1.000 & 0.807 & 0.512 & 0.280 & 0.138 \\
  $|\overline{S}|$  & 0.445 & 0.689 & 0.938 & 1.000 & 0.813 & 0.520 & 0.285 & 0.141 \\
  $|\Delta\ln u|$   & 0.997 & 0.995 & 0.993 & 0.991 & 0.990 & 0.992 & 0.996 & 1.000 \\
  \bottomrule
\end{tabular}
\end{table}

The first two rows agree to $2.2\,\%$ at every column; the third is flat to one
per cent. The density structure of the plateau error is the density structure of
the observable's sensitivity to its reservoir.

That sensitivity has an interior maximum for the reason
Section~\ref{sec:ground_fed_fraction} gives. Both $\ffrac{3}$ and $\ffrac{4}$
rise from zero to unity as the reservoir grows, but the two shells switch from
recombination-fed to ground-fed at different reservoir strengths, because the
ground-fed channel weakens with principal quantum number faster than the
recombination-fed one. Their difference therefore vanishes in both supply limits
and peaks between them. The density ridge is where the reservoir sits in that
mixed-supply region.

\subsection{Alternatives that would produce the same appearance}
\label{sec:ridge_alternatives}

An interior maximum on a coarse grid can be manufactured in several ways, and
the attribution above does not by itself exclude them.

\emph{Grid resolution.} The ridge occupies one density column and the grid is
spaced by a factor $2.683$. Its location is therefore known only to that factor,
and this work cannot say whether the true maximum lies at
$1.93\times10^{13}\,\si{\per\cubic\centi\metre}$ or anywhere within a factor
$2.7$ of it. What is established is that the maximum is interior --- not at
either end of the density range --- and that it does not move between
temperature rows. Refining the density grid requires rebuilding the rate matrix
and revalidating every number in Chapter~\ref{ch:theory} against the new one;
future work could determine the location more precisely.

\emph{Step-size convention.} A ridge whose position depended on the size of the
imposed step would be an artifact of that choice. It does not: the error maximum
falls on the same density column for a one-interval and a four-interval step in
$31$ of the $31$ temperature rows both can reach, and so does the maximum of
$|\overline{S}|$ (\texttt{verify\_ridge\_mechanism.py}). This tests steps of
fixed \emph{fractional} size; a step of fixed absolute size in electron-volts
spans a different number of grid intervals at each temperature and is a
different experiment, not tested here.

\emph{Truncation.} The model resolves $n\le8$ and bundles $n=9$--$15$. The
mechanism turns on the difference between the $n=3$ and $n=4$ supply fractions,
both fully resolved, but those fractions depend on the recombination flux
cascading from above, which the bundling approximates. This work does not test
the sensitivity of the ridge to $n_{\max}$; Chapter~\ref{ch:limits} lists it
among the outstanding checks.

\subsection{Across temperature: a displacement effect}
\label{sec:temperature_mechanism}

The temperature dependence has the opposite explanation, and this is where
treating the two axes alike would go wrong.

Along the ridge column the plateau error falls monotonically from $18.07\,\%$ at
$\Te = 2.02\,$eV to $4.55\,\%$ at $8.29\,$eV, a factor $3.72$. Decomposed
through Eq.~\eqref{eq:exact_decomposition}, $|\Delta\ln u|$ contributes a factor
$2.72$ of that fall and $|\overline{S}|$ only $1.29$; in logarithmic terms the
displacement carries $76\,\%$ of the decay and the sensitivity $19\,\%$. At the
four-interval step the split is $72\,\%$ and $13\,\%$
(\texttt{verify\_ridge\_mechanism.py}).

The observable therefore does not become much less sensitive as the plasma
heats. \textbf{The ground state simply moves less.} A step of fixed fractional
size displaces the reservoir progressively less at higher temperature, and the
error follows it down.

Why the displacement behaves this way is not established here. In particular
$|\Delta\ln u|$ is flat to one per cent across three decades of density while
varying by a factor $2.7$ across temperature, which implies that
$\partial\ln\ngs/\partial\ln\Te$ at collisional--radiative equilibrium is
controlled by temperature alone. \emph{[MECHANISM NOT ESTABLISHED: this is an
observation from \texttt{verify\_ridge\_mechanism.py}, not a derived result. It
has not been shown analytically and is a candidate for future work.]}

There is no interior maximum in temperature anywhere in the defended range. The
largest values, near $\Te\approx2\,$eV, are boundary values of the
optical-thickness restriction of Section~\ref{sec:optical_depth}, not a peak.

\subsection{Three explanations this work does not support}
\label{sec:withdrawn}

Three readings of the ridge are natural and do not survive the numbers. Each is
recorded because it is among the first a reader is likely to propose.

\textbf{It is not detachment.} The ridge sits at
$1.93\times10^{19}\,\si{\per\cubic\metre}$: a factor $5.2$ below the density
band that published edge modelling associates with a detached ITER divertor
\emph{[SOURCE REQUIRED: Guillemaut et al.\ 2011, Fus.\ Eng.\ Des.\ 86, 2954 ---
key not yet in the bibliography]}, and a factor $1.6$ below the modelled
upstream separatrix density \emph{[SOURCE REQUIRED: Stangeby et al.\ 2022,
Nucl.\ Fusion 62]}. It falls between the two, in a region neither source
characterises. The ridge is reproducible; its identification with a named
divertor operating regime is not supported by anything in this work.

\textbf{It is not a peak in temperature.} The error decreases monotonically with
temperature throughout the defended range, as
Section~\ref{sec:temperature_mechanism} shows.

\textbf{It is not simply where $\bdep^{\rm CRE}$ meets the sensitivity
maximum.} That reading accounts for the density structure but fails on the
temperature axis: at the four-interval step the ratio
$\bdep^{\rm CRE}/\bdep^{\rm peak}$ passes through unity at $\Te = 4.7\,$eV,
where the error is $33\,\%$ and falling monotonically
(\texttt{verify\_plateau\_gridmap.py}). The endpoint of the reservoir
trajectory is the wrong object: Eq.~\eqref{eq:Sbar_def} shows that what enters
is the mean sensitivity over the whole interval traversed.

\subsection{What this section establishes}

The plateau error factorises exactly into a sensitivity and a displacement,
Eq.~\eqref{eq:eps_plateau_exact}. Attributing the structure of the error maps to
one factor or the other gives different answers on the two axes: across density
the sensitivity carries the structure while the displacement is flat to one per
cent, and across temperature the displacement carries three quarters of the
decay. The ridge is a property of how the $n=3$ and $n=4$ shells divide their
supply between excitation from the ground state and recombination from the
continuum --- a property of the atomic system, needing no divertor geometry to
produce it --- while its amplitude at any density is set by how far a
temperature step moves the ground state.

That the local model generates a ridge without spatial structure does not imply
spatial structure is irrelevant to a real line-integrated measurement, which
Chapter~\ref{ch:limits} takes up.

The argument so far is entirely algebraic. Every quantity above is built from
the partial-equilibrium state $\Robs^{\rm PE}$, and that state was constructed
by freezing the ground state by hand rather than by integrating anything.
Whether the time-dependent system actually visits it is a separate question.


% ------------------------------------------------------------------
% Sections still to be written. The labels are declared here so that the
% cross-references in chapter5_C5D.tex resolve rather than printing "??".
% Replace each stub with the real section as it is drafted.
% ------------------------------------------------------------------
\section{Does the transient visit the plateau?}
\label{sec:bridge}

\emph{[TO BE WRITTEN --- C5-A writeup. Full 43-state integration; the
dynamic bridge; the closed form for the early-response slope $k$, which was
cut from Section~\ref{sec:mechanism} because it needs the transient
established first.]}

\section{How long the error persists}
\label{sec:persistence}

\emph{[TO BE WRITTEN. $\tau_{\rm eff}/\tauqss = 1.01$ and $1.07$.]}

\section{Optical thickness and the restricted range}
\label{sec:optical_depth}

\emph{[TO BE WRITTEN --- may instead live in Chapter~\ref{ch:theory};
Section~\ref{sec:mechanism} refers to it for the $\Te \ge 2$~eV restriction.]}
          % if the other chapters use \include
%   or
%     % ============================================================
% chapter5.tex -- Chapter 5 wrapper.
%
% chapter5_C5D.tex is a SECTION (\section{...}), not a chapter. It needs a
% \chapter{} above it or LaTeX has nowhere to attach it.
%
% Add to thesis_main.tex, in the document body, after chapter 4:
%
%     % ============================================================
% chapter5.tex -- Chapter 5 wrapper.
%
% chapter5_C5D.tex is a SECTION (\section{...}), not a chapter. It needs a
% \chapter{} above it or LaTeX has nowhere to attach it.
%
% Add to thesis_main.tex, in the document body, after chapter 4:
%
%     \include{chapter5}          % if the other chapters use \include
%   or
%     \input{chapter5}            % if they use \input
%
% Use whichever the neighbouring chapters use -- mixing them causes the page
% breaks and numbering to differ from the rest of the thesis.
% ============================================================

\chapter{What the Model Says}
\label{ch:results}

% ------------------------------------------------------------------
% PLACEHOLDER so the chapter is never empty while sections are drafted.
% Delete this block once sec:error_maps exists.
% ------------------------------------------------------------------
\section{The error maps}
\label{sec:error_maps}

\emph{[TO BE WRITTEN --- C5-F. Defines the operating grid, presents
$\epsstep$ and $\epsplat$ over it, and gives the headline numbers with their
scopes. Section~\ref{sec:mechanism} below is drafted and refers back to this
section for the maps it explains.]}

% ------------------------------------------------------------------
% C5-D, drafted 24 August 2026.
% ------------------------------------------------------------------
\input{chapter5_C5D}

% ------------------------------------------------------------------
% Sections still to be written. The labels are declared here so that the
% cross-references in chapter5_C5D.tex resolve rather than printing "??".
% Replace each stub with the real section as it is drafted.
% ------------------------------------------------------------------
\section{Does the transient visit the plateau?}
\label{sec:bridge}

\emph{[TO BE WRITTEN --- C5-A writeup. Full 43-state integration; the
dynamic bridge; the closed form for the early-response slope $k$, which was
cut from Section~\ref{sec:mechanism} because it needs the transient
established first.]}

\section{How long the error persists}
\label{sec:persistence}

\emph{[TO BE WRITTEN. $\tau_{\rm eff}/\tauqss = 1.01$ and $1.07$.]}

\section{Optical thickness and the restricted range}
\label{sec:optical_depth}

\emph{[TO BE WRITTEN --- may instead live in Chapter~\ref{ch:theory};
Section~\ref{sec:mechanism} refers to it for the $\Te \ge 2$~eV restriction.]}
          % if the other chapters use \include
%   or
%     % ============================================================
% chapter5.tex -- Chapter 5 wrapper.
%
% chapter5_C5D.tex is a SECTION (\section{...}), not a chapter. It needs a
% \chapter{} above it or LaTeX has nowhere to attach it.
%
% Add to thesis_main.tex, in the document body, after chapter 4:
%
%     \include{chapter5}          % if the other chapters use \include
%   or
%     \input{chapter5}            % if they use \input
%
% Use whichever the neighbouring chapters use -- mixing them causes the page
% breaks and numbering to differ from the rest of the thesis.
% ============================================================

\chapter{What the Model Says}
\label{ch:results}

% ------------------------------------------------------------------
% PLACEHOLDER so the chapter is never empty while sections are drafted.
% Delete this block once sec:error_maps exists.
% ------------------------------------------------------------------
\section{The error maps}
\label{sec:error_maps}

\emph{[TO BE WRITTEN --- C5-F. Defines the operating grid, presents
$\epsstep$ and $\epsplat$ over it, and gives the headline numbers with their
scopes. Section~\ref{sec:mechanism} below is drafted and refers back to this
section for the maps it explains.]}

% ------------------------------------------------------------------
% C5-D, drafted 24 August 2026.
% ------------------------------------------------------------------
\input{chapter5_C5D}

% ------------------------------------------------------------------
% Sections still to be written. The labels are declared here so that the
% cross-references in chapter5_C5D.tex resolve rather than printing "??".
% Replace each stub with the real section as it is drafted.
% ------------------------------------------------------------------
\section{Does the transient visit the plateau?}
\label{sec:bridge}

\emph{[TO BE WRITTEN --- C5-A writeup. Full 43-state integration; the
dynamic bridge; the closed form for the early-response slope $k$, which was
cut from Section~\ref{sec:mechanism} because it needs the transient
established first.]}

\section{How long the error persists}
\label{sec:persistence}

\emph{[TO BE WRITTEN. $\tau_{\rm eff}/\tauqss = 1.01$ and $1.07$.]}

\section{Optical thickness and the restricted range}
\label{sec:optical_depth}

\emph{[TO BE WRITTEN --- may instead live in Chapter~\ref{ch:theory};
Section~\ref{sec:mechanism} refers to it for the $\Te \ge 2$~eV restriction.]}
            % if they use \input
%
% Use whichever the neighbouring chapters use -- mixing them causes the page
% breaks and numbering to differ from the rest of the thesis.
% ============================================================

\chapter{What the Model Says}
\label{ch:results}

% ------------------------------------------------------------------
% PLACEHOLDER so the chapter is never empty while sections are drafted.
% Delete this block once sec:error_maps exists.
% ------------------------------------------------------------------
\section{The error maps}
\label{sec:error_maps}

\emph{[TO BE WRITTEN --- C5-F. Defines the operating grid, presents
$\epsstep$ and $\epsplat$ over it, and gives the headline numbers with their
scopes. Section~\ref{sec:mechanism} below is drafted and refers back to this
section for the maps it explains.]}

% ------------------------------------------------------------------
% C5-D, drafted 24 August 2026.
% ------------------------------------------------------------------
% ============================================================
% Chapter 5, Section: the mechanism.  C5-D, draft 2.
%
% Written under the thesis-writing skill, 24 August 2026.
% Source: CH5_EVIDENCE.md. L_grid SHA 2d92b58e..., S_grid SHA 7822f536...
% Draft 1 retained as chapter5_C5D_draft1_20260824.tex (failed six gates).
%
% MACROS beyond those chapter3.tex exercises:  \epsstep  \epsplat
%   (recorded in notation_and_definitions.md Sec. 7, never yet compiled --
%    CONFIRM they exist in thesis_main.tex before building.)
% ============================================================

%-----------------------------------------------------------------------------
\section{Why the error is largest where it is}
\label{sec:mechanism}
%-----------------------------------------------------------------------------

The maps of Section~\ref{sec:error_maps} show where the diagnostic error is
large. They do not say why, and a map without a mechanism is a lookup table
that cannot be extrapolated past the grid that produced it. This section asks
which physical quantity puts the ridge where it is.

The answer has two halves and they disagree. The structure along the density
axis and the structure along the temperature axis have different causes, and
reading them as one effect --- which the appearance of a single ridged surface
invites --- gets both wrong.

\subsection{Two error measures}
\label{sec:eps_definitions}

Section~\ref{sec:two_references} fixed the two reference states against which an
observable can be judged. Two further quantities are needed here, and both are
properties of a temperature \emph{step} rather than of a single operator, which
is why neither appears in Chapter~\ref{ch:theory}.

Let the plasma sit at the collisional--radiative equilibrium of
$\Lmat(\Te,\nel)$, and let the electron temperature step to $\Te+\Delta\Te$ at
$t=0$ with $\nel$ held fixed. Write $\Lmat^-$ and $\Lmat^+$ for the operators
before and after, in the notation of Section~\ref{sec:two_clocks}. The
\emph{step error}
\begin{equation}
  \epsstep = \left|\frac{\Rcre(\Te,\nel)}
                        {\Rcre(\Te+\Delta\Te,\nel)} - 1\right|
  \label{eq:eps_step_def}
\end{equation}

\noindent is the dimensionless distance between the two tabulated answers: what
the reading would change by if the plasma re-equilibrated instantly. The
\emph{plateau error}
\begin{equation}
  \epsplat = \left|\frac{\Robs^{\rm PE}}
                        {\Rcre(\Te+\Delta\Te,\nel)} - 1\right|,
  \qquad
  \Robs^{\rm PE} \equiv \Rqss\!\left(\Te+\Delta\Te,\nel;\,\bdep^-\right),
  \label{eq:eps_plateau_def}
\end{equation}

\noindent is the distance between the tabulated answer and the state the plasma
occupies once the excited manifold has settled but before the ground state has
moved. The superscript PE denotes \emph{partial equilibrium}: the post-step
excited response of Eq.~\eqref{eq:qss_ratio}, evaluated on the \emph{pre-step}
reservoir.

The consequence of keeping these separate is that they are not close. Across the
operating grid $\epsplat$ exceeds $\epsstep$ at every point tested, typically by
an order of magnitude (Section~\ref{sec:error_maps}). A diagnostician who
pictures the reading moving smoothly between the two tabulated values has not
merely misjudged the size of the excursion but has the wrong shape of transient.
$\epsplat$ is the quantity this chapter reports, and Section~\ref{sec:bridge}
establishes that the time-dependent system does pass through the state it
measures.

\subsection{An exact decomposition of the finite step}
\label{sec:exact_decomposition}

Chapter~\ref{ch:theory} worked in the ground-state departure coefficient
$\bdep = \ngs/(\Zsb\nel\nion)$, because that is the variable in which the
observable's response has the closed form of Eq.~\eqref{eq:R_of_b}. Here it is
better to use the reservoir itself,
\begin{equation}
  u \;\equiv\; \frac{\ngs}{\nion} \;=\; \bdep\,\Zsb(\Te)\,\nel ,
  \label{eq:u_def}
\end{equation}

\noindent a dimensionless ground-state population per unit ion density. At
fixed $(\Te,\nel)$ the two are proportional, so $\mathrm{d}\ln u =
\mathrm{d}\ln\bdep$ and the sensitivity of
Section~\ref{sec:ground_fed_fraction} carries over unchanged. The reason for
the change of variable is that a temperature step moves $\Zsb$: comparing a
pre-step $\bdep$ with a post-step one would mix a change in the reservoir with
a change in the Saha--Boltzmann normalisation, and $u$ has no normalisation to
move. That confusion is not hypothetical --- an earlier version of the analysis
reported in this section made exactly it, and the resulting contamination in
$\Delta\ln\bdep$ was comparable to the quantity being measured.

\textbf{Assumption.} Equation~\eqref{eq:u_def} treats $\nion$ as a fixed
reservoir, unchanged by the step, as throughout this thesis
(Section~\ref{sec:rate_equation}). This removes the freedom for the ionisation
balance to respond, and it is what makes $u$ numerically equal to the
ground-state population in the solved system. Chapter~\ref{ch:limits} returns
to what a mobile ion population would change.

With that, Eq.~\eqref{eq:R_of_b} makes the post-step observable a ratio of two
affine functions of $u$, so the logarithmic sensitivity of
Section~\ref{sec:ground_fed_fraction} is $\mathrm{d}\ln\Robs/\mathrm{d}\ln u =
\ffrac{3}-\ffrac{4}$, the difference of the two ground-fed fractions.
Integrating between the frozen and the equilibrated reservoir gives, with no
approximation,
\begin{equation}
  \ln\frac{\Robs^{\rm PE}}{\Rcre}
  \;=\; \int_{\ln u^+}^{\ln u^-}\!\!\big(\ffrac{3}-\ffrac{4}\big)\,\mathrm{d}\ln u
  \;=\; \overline{S}\;\Delta\ln u ,
  \label{eq:exact_decomposition}
\end{equation}

\noindent where $\Delta\ln u = \ln(u^-/u^+)$ is the reservoir displacement and
\begin{equation}
  \overline{S} \;\equiv\;
  \frac{1}{\Delta\ln u}\int_{\ln u^+}^{\ln u^-}\!\!
  \big(\ffrac{3}-\ffrac{4}\big)\,\mathrm{d}\ln u
  \label{eq:Sbar_def}
\end{equation}

\noindent is the mean sensitivity over the interval the reservoir traverses.
Both factors are dimensionless. Since $\ffrac{3},\ffrac{4}\in[0,1]$ by
Section~\ref{sec:ground_fed_fraction}, and $\overline{S}$ is an average of their
difference, $|\overline{S}|$ is bounded by the peak sensitivity of
Eq.~\eqref{eq:peak_sensitivity} and so is strictly less than unity everywhere.
The plateau error follows:
\begin{equation}
  \boxed{\ \epsplat = \left|\,\mathrm{e}^{\overline{S}\,\Delta\ln u} - 1\,\right|\ }
  \label{eq:eps_plateau_exact}
\end{equation}

The consequence worth extracting is that a single scalar error has become a
product of two quantities that can be measured independently: how strongly the
observable responds to its reservoir, and how far the reservoir moved. Nothing
in Eq.~\eqref{eq:exact_decomposition} says which of them varies across the
operating map, and that is what the rest of this section determines.

\textbf{Use Eq.~\eqref{eq:eps_plateau_exact}, not its linearisation.} The
first-order form $\epsplat \approx |\overline{S}\,\Delta\ln u|$ is adequate for
a small step and wrong for a large one: at the four-interval step used below it
returns $59.4\,\%$ where the exact expression gives $81.1\,\%$, understating by
more than a quarter (\texttt{verify\_ridge\_mechanism.py}). The exponential
costs nothing.

One caution about what follows. Because Eq.~\eqref{eq:exact_decomposition} is an
identity, comparing its two sides tests nothing: a check that $\epsplat$ equals
$|\mathrm{e}^{\overline{S}\Delta\ln u}-1|$ can fail only through arithmetic
error. What the identity permits is an \emph{attribution} --- the two factors
are separately measurable, so one may ask which carries the structure the maps
show. That is a weaker kind of statement than a derivation, and it is the kind
made below.

\subsection{Across density: a sensitivity structure}
\label{sec:density_mechanism}

The plateau error has an interior maximum in density at
$\nel \approx 1.93\times10^{13}\,\si{\per\cubic\centi\metre}$, at every
temperature in the range this thesis defends.

Evaluating both factors of Eq.~\eqref{eq:exact_decomposition} at every grid
point settles which of them puts it there
(\texttt{verify\_ridge\_mechanism.py}, over $\Te = 2$--$9\,$eV and all eight
density columns). The column maximising $\epsplat$ is the column maximising
$|\overline{S}|$ in $32$ of $32$ temperature rows for a one-interval
temperature step and $31$ of $31$ for a four-interval step; the column
maximising $|\Delta\ln u|$ coincides with the error maximum in $0$ of $32$ and
$0$ of $31$.

Coincidence counts are weak evidence on their own, because a quantity that
barely varies still has its maximum somewhere. The amplitudes settle it. Across
the three decades of density $|\overline{S}|$ varies by a factor $7.11$ and
$|\Delta\ln u|$ by a factor $1.01$. \textbf{A quantity that changes by one per
cent cannot produce a ridge in which the error changes by a factor of seven},
wherever its own maximum falls. Table~\ref{tab:mechanism_profiles} makes the
point without statistics.

\begin{table}[htbp]
\centering
\caption{Density profiles of the plateau error and of the two factors of
Eq.~\eqref{eq:exact_decomposition}, each normalised to its own row maximum and
then taken as the median over the $32$ temperature rows with
$\Te = 2$--$9\,$eV, for a one-interval temperature step. From
\texttt{verify\_ridge\_mechanism.py}.}
\label{tab:mechanism_profiles}
\sisetup{table-format=1.3}
\begin{tabular}{lSSSSSSSS}
  \toprule
  $\nel$ (\si{\per\cubic\centi\metre}) & {$10^{12}$} & {$2.7{\cdot}10^{12}$} &
  {$7.2{\cdot}10^{12}$} & {$1.9{\cdot}10^{13}$} & {$5.2{\cdot}10^{13}$} &
  {$1.4{\cdot}10^{14}$} & {$3.7{\cdot}10^{14}$} & {$10^{15}$} \\
  \midrule
  $\epsplat$        & 0.437 & 0.683 & 0.937 & 1.000 & 0.807 & 0.512 & 0.280 & 0.138 \\
  $|\overline{S}|$  & 0.445 & 0.689 & 0.938 & 1.000 & 0.813 & 0.520 & 0.285 & 0.141 \\
  $|\Delta\ln u|$   & 0.997 & 0.995 & 0.993 & 0.991 & 0.990 & 0.992 & 0.996 & 1.000 \\
  \bottomrule
\end{tabular}
\end{table}

The first two rows agree to $2.2\,\%$ at every column; the third is flat to one
per cent. The density structure of the plateau error is the density structure of
the observable's sensitivity to its reservoir.

That sensitivity has an interior maximum for the reason
Section~\ref{sec:ground_fed_fraction} gives. Both $\ffrac{3}$ and $\ffrac{4}$
rise from zero to unity as the reservoir grows, but the two shells switch from
recombination-fed to ground-fed at different reservoir strengths, because the
ground-fed channel weakens with principal quantum number faster than the
recombination-fed one. Their difference therefore vanishes in both supply limits
and peaks between them. The density ridge is where the reservoir sits in that
mixed-supply region.

\subsection{Alternatives that would produce the same appearance}
\label{sec:ridge_alternatives}

An interior maximum on a coarse grid can be manufactured in several ways, and
the attribution above does not by itself exclude them.

\emph{Grid resolution.} The ridge occupies one density column and the grid is
spaced by a factor $2.683$. Its location is therefore known only to that factor,
and this work cannot say whether the true maximum lies at
$1.93\times10^{13}\,\si{\per\cubic\centi\metre}$ or anywhere within a factor
$2.7$ of it. What is established is that the maximum is interior --- not at
either end of the density range --- and that it does not move between
temperature rows. Refining the density grid requires rebuilding the rate matrix
and revalidating every number in Chapter~\ref{ch:theory} against the new one;
future work could determine the location more precisely.

\emph{Step-size convention.} A ridge whose position depended on the size of the
imposed step would be an artifact of that choice. It does not: the error maximum
falls on the same density column for a one-interval and a four-interval step in
$31$ of the $31$ temperature rows both can reach, and so does the maximum of
$|\overline{S}|$ (\texttt{verify\_ridge\_mechanism.py}). This tests steps of
fixed \emph{fractional} size; a step of fixed absolute size in electron-volts
spans a different number of grid intervals at each temperature and is a
different experiment, not tested here.

\emph{Truncation.} The model resolves $n\le8$ and bundles $n=9$--$15$. The
mechanism turns on the difference between the $n=3$ and $n=4$ supply fractions,
both fully resolved, but those fractions depend on the recombination flux
cascading from above, which the bundling approximates. This work does not test
the sensitivity of the ridge to $n_{\max}$; Chapter~\ref{ch:limits} lists it
among the outstanding checks.

\subsection{Across temperature: a displacement effect}
\label{sec:temperature_mechanism}

The temperature dependence has the opposite explanation, and this is where
treating the two axes alike would go wrong.

Along the ridge column the plateau error falls monotonically from $18.07\,\%$ at
$\Te = 2.02\,$eV to $4.55\,\%$ at $8.29\,$eV, a factor $3.72$. Decomposed
through Eq.~\eqref{eq:exact_decomposition}, $|\Delta\ln u|$ contributes a factor
$2.72$ of that fall and $|\overline{S}|$ only $1.29$; in logarithmic terms the
displacement carries $76\,\%$ of the decay and the sensitivity $19\,\%$. At the
four-interval step the split is $72\,\%$ and $13\,\%$
(\texttt{verify\_ridge\_mechanism.py}).

The observable therefore does not become much less sensitive as the plasma
heats. \textbf{The ground state simply moves less.} A step of fixed fractional
size displaces the reservoir progressively less at higher temperature, and the
error follows it down.

Why the displacement behaves this way is not established here. In particular
$|\Delta\ln u|$ is flat to one per cent across three decades of density while
varying by a factor $2.7$ across temperature, which implies that
$\partial\ln\ngs/\partial\ln\Te$ at collisional--radiative equilibrium is
controlled by temperature alone. \emph{[MECHANISM NOT ESTABLISHED: this is an
observation from \texttt{verify\_ridge\_mechanism.py}, not a derived result. It
has not been shown analytically and is a candidate for future work.]}

There is no interior maximum in temperature anywhere in the defended range. The
largest values, near $\Te\approx2\,$eV, are boundary values of the
optical-thickness restriction of Section~\ref{sec:optical_depth}, not a peak.

\subsection{Three explanations this work does not support}
\label{sec:withdrawn}

Three readings of the ridge are natural and do not survive the numbers. Each is
recorded because it is among the first a reader is likely to propose.

\textbf{It is not detachment.} The ridge sits at
$1.93\times10^{19}\,\si{\per\cubic\metre}$: a factor $5.2$ below the density
band that published edge modelling associates with a detached ITER divertor
\emph{[SOURCE REQUIRED: Guillemaut et al.\ 2011, Fus.\ Eng.\ Des.\ 86, 2954 ---
key not yet in the bibliography]}, and a factor $1.6$ below the modelled
upstream separatrix density \emph{[SOURCE REQUIRED: Stangeby et al.\ 2022,
Nucl.\ Fusion 62]}. It falls between the two, in a region neither source
characterises. The ridge is reproducible; its identification with a named
divertor operating regime is not supported by anything in this work.

\textbf{It is not a peak in temperature.} The error decreases monotonically with
temperature throughout the defended range, as
Section~\ref{sec:temperature_mechanism} shows.

\textbf{It is not simply where $\bdep^{\rm CRE}$ meets the sensitivity
maximum.} That reading accounts for the density structure but fails on the
temperature axis: at the four-interval step the ratio
$\bdep^{\rm CRE}/\bdep^{\rm peak}$ passes through unity at $\Te = 4.7\,$eV,
where the error is $33\,\%$ and falling monotonically
(\texttt{verify\_plateau\_gridmap.py}). The endpoint of the reservoir
trajectory is the wrong object: Eq.~\eqref{eq:Sbar_def} shows that what enters
is the mean sensitivity over the whole interval traversed.

\subsection{What this section establishes}

The plateau error factorises exactly into a sensitivity and a displacement,
Eq.~\eqref{eq:eps_plateau_exact}. Attributing the structure of the error maps to
one factor or the other gives different answers on the two axes: across density
the sensitivity carries the structure while the displacement is flat to one per
cent, and across temperature the displacement carries three quarters of the
decay. The ridge is a property of how the $n=3$ and $n=4$ shells divide their
supply between excitation from the ground state and recombination from the
continuum --- a property of the atomic system, needing no divertor geometry to
produce it --- while its amplitude at any density is set by how far a
temperature step moves the ground state.

That the local model generates a ridge without spatial structure does not imply
spatial structure is irrelevant to a real line-integrated measurement, which
Chapter~\ref{ch:limits} takes up.

The argument so far is entirely algebraic. Every quantity above is built from
the partial-equilibrium state $\Robs^{\rm PE}$, and that state was constructed
by freezing the ground state by hand rather than by integrating anything.
Whether the time-dependent system actually visits it is a separate question.


% ------------------------------------------------------------------
% Sections still to be written. The labels are declared here so that the
% cross-references in chapter5_C5D.tex resolve rather than printing "??".
% Replace each stub with the real section as it is drafted.
% ------------------------------------------------------------------
\section{Does the transient visit the plateau?}
\label{sec:bridge}

\emph{[TO BE WRITTEN --- C5-A writeup. Full 43-state integration; the
dynamic bridge; the closed form for the early-response slope $k$, which was
cut from Section~\ref{sec:mechanism} because it needs the transient
established first.]}

\section{How long the error persists}
\label{sec:persistence}

\emph{[TO BE WRITTEN. $\tau_{\rm eff}/\tauqss = 1.01$ and $1.07$.]}

\section{Optical thickness and the restricted range}
\label{sec:optical_depth}

\emph{[TO BE WRITTEN --- may instead live in Chapter~\ref{ch:theory};
Section~\ref{sec:mechanism} refers to it for the $\Te \ge 2$~eV restriction.]}
            % if they use \input
%
% Use whichever the neighbouring chapters use -- mixing them causes the page
% breaks and numbering to differ from the rest of the thesis.
% ============================================================

\chapter{What the Model Says}
\label{ch:results}

% ------------------------------------------------------------------
% PLACEHOLDER so the chapter is never empty while sections are drafted.
% Delete this block once sec:error_maps exists.
% ------------------------------------------------------------------
\section{The error maps}
\label{sec:error_maps}

\emph{[TO BE WRITTEN --- C5-F. Defines the operating grid, presents
$\epsstep$ and $\epsplat$ over it, and gives the headline numbers with their
scopes. Section~\ref{sec:mechanism} below is drafted and refers back to this
section for the maps it explains.]}

% ------------------------------------------------------------------
% C5-D, drafted 24 August 2026.
% ------------------------------------------------------------------
% ============================================================
% Chapter 5, Section: the mechanism.  C5-D, draft 2.
%
% Written under the thesis-writing skill, 24 August 2026.
% Source: CH5_EVIDENCE.md. L_grid SHA 2d92b58e..., S_grid SHA 7822f536...
% Draft 1 retained as chapter5_C5D_draft1_20260824.tex (failed six gates).
%
% MACROS beyond those chapter3.tex exercises:  \epsstep  \epsplat
%   (recorded in notation_and_definitions.md Sec. 7, never yet compiled --
%    CONFIRM they exist in thesis_main.tex before building.)
% ============================================================

%-----------------------------------------------------------------------------
\section{Why the error is largest where it is}
\label{sec:mechanism}
%-----------------------------------------------------------------------------

The maps of Section~\ref{sec:error_maps} show where the diagnostic error is
large. They do not say why, and a map without a mechanism is a lookup table
that cannot be extrapolated past the grid that produced it. This section asks
which physical quantity puts the ridge where it is.

The answer has two halves and they disagree. The structure along the density
axis and the structure along the temperature axis have different causes, and
reading them as one effect --- which the appearance of a single ridged surface
invites --- gets both wrong.

\subsection{Two error measures}
\label{sec:eps_definitions}

Section~\ref{sec:two_references} fixed the two reference states against which an
observable can be judged. Two further quantities are needed here, and both are
properties of a temperature \emph{step} rather than of a single operator, which
is why neither appears in Chapter~\ref{ch:theory}.

Let the plasma sit at the collisional--radiative equilibrium of
$\Lmat(\Te,\nel)$, and let the electron temperature step to $\Te+\Delta\Te$ at
$t=0$ with $\nel$ held fixed. Write $\Lmat^-$ and $\Lmat^+$ for the operators
before and after, in the notation of Section~\ref{sec:two_clocks}. The
\emph{step error}
\begin{equation}
  \epsstep = \left|\frac{\Rcre(\Te,\nel)}
                        {\Rcre(\Te+\Delta\Te,\nel)} - 1\right|
  \label{eq:eps_step_def}
\end{equation}

\noindent is the dimensionless distance between the two tabulated answers: what
the reading would change by if the plasma re-equilibrated instantly. The
\emph{plateau error}
\begin{equation}
  \epsplat = \left|\frac{\Robs^{\rm PE}}
                        {\Rcre(\Te+\Delta\Te,\nel)} - 1\right|,
  \qquad
  \Robs^{\rm PE} \equiv \Rqss\!\left(\Te+\Delta\Te,\nel;\,\bdep^-\right),
  \label{eq:eps_plateau_def}
\end{equation}

\noindent is the distance between the tabulated answer and the state the plasma
occupies once the excited manifold has settled but before the ground state has
moved. The superscript PE denotes \emph{partial equilibrium}: the post-step
excited response of Eq.~\eqref{eq:qss_ratio}, evaluated on the \emph{pre-step}
reservoir.

The consequence of keeping these separate is that they are not close. Across the
operating grid $\epsplat$ exceeds $\epsstep$ at every point tested, typically by
an order of magnitude (Section~\ref{sec:error_maps}). A diagnostician who
pictures the reading moving smoothly between the two tabulated values has not
merely misjudged the size of the excursion but has the wrong shape of transient.
$\epsplat$ is the quantity this chapter reports, and Section~\ref{sec:bridge}
establishes that the time-dependent system does pass through the state it
measures.

\subsection{An exact decomposition of the finite step}
\label{sec:exact_decomposition}

Chapter~\ref{ch:theory} worked in the ground-state departure coefficient
$\bdep = \ngs/(\Zsb\nel\nion)$, because that is the variable in which the
observable's response has the closed form of Eq.~\eqref{eq:R_of_b}. Here it is
better to use the reservoir itself,
\begin{equation}
  u \;\equiv\; \frac{\ngs}{\nion} \;=\; \bdep\,\Zsb(\Te)\,\nel ,
  \label{eq:u_def}
\end{equation}

\noindent a dimensionless ground-state population per unit ion density. At
fixed $(\Te,\nel)$ the two are proportional, so $\mathrm{d}\ln u =
\mathrm{d}\ln\bdep$ and the sensitivity of
Section~\ref{sec:ground_fed_fraction} carries over unchanged. The reason for
the change of variable is that a temperature step moves $\Zsb$: comparing a
pre-step $\bdep$ with a post-step one would mix a change in the reservoir with
a change in the Saha--Boltzmann normalisation, and $u$ has no normalisation to
move. That confusion is not hypothetical --- an earlier version of the analysis
reported in this section made exactly it, and the resulting contamination in
$\Delta\ln\bdep$ was comparable to the quantity being measured.

\textbf{Assumption.} Equation~\eqref{eq:u_def} treats $\nion$ as a fixed
reservoir, unchanged by the step, as throughout this thesis
(Section~\ref{sec:rate_equation}). This removes the freedom for the ionisation
balance to respond, and it is what makes $u$ numerically equal to the
ground-state population in the solved system. Chapter~\ref{ch:limits} returns
to what a mobile ion population would change.

With that, Eq.~\eqref{eq:R_of_b} makes the post-step observable a ratio of two
affine functions of $u$, so the logarithmic sensitivity of
Section~\ref{sec:ground_fed_fraction} is $\mathrm{d}\ln\Robs/\mathrm{d}\ln u =
\ffrac{3}-\ffrac{4}$, the difference of the two ground-fed fractions.
Integrating between the frozen and the equilibrated reservoir gives, with no
approximation,
\begin{equation}
  \ln\frac{\Robs^{\rm PE}}{\Rcre}
  \;=\; \int_{\ln u^+}^{\ln u^-}\!\!\big(\ffrac{3}-\ffrac{4}\big)\,\mathrm{d}\ln u
  \;=\; \overline{S}\;\Delta\ln u ,
  \label{eq:exact_decomposition}
\end{equation}

\noindent where $\Delta\ln u = \ln(u^-/u^+)$ is the reservoir displacement and
\begin{equation}
  \overline{S} \;\equiv\;
  \frac{1}{\Delta\ln u}\int_{\ln u^+}^{\ln u^-}\!\!
  \big(\ffrac{3}-\ffrac{4}\big)\,\mathrm{d}\ln u
  \label{eq:Sbar_def}
\end{equation}

\noindent is the mean sensitivity over the interval the reservoir traverses.
Both factors are dimensionless. Since $\ffrac{3},\ffrac{4}\in[0,1]$ by
Section~\ref{sec:ground_fed_fraction}, and $\overline{S}$ is an average of their
difference, $|\overline{S}|$ is bounded by the peak sensitivity of
Eq.~\eqref{eq:peak_sensitivity} and so is strictly less than unity everywhere.
The plateau error follows:
\begin{equation}
  \boxed{\ \epsplat = \left|\,\mathrm{e}^{\overline{S}\,\Delta\ln u} - 1\,\right|\ }
  \label{eq:eps_plateau_exact}
\end{equation}

The consequence worth extracting is that a single scalar error has become a
product of two quantities that can be measured independently: how strongly the
observable responds to its reservoir, and how far the reservoir moved. Nothing
in Eq.~\eqref{eq:exact_decomposition} says which of them varies across the
operating map, and that is what the rest of this section determines.

\textbf{Use Eq.~\eqref{eq:eps_plateau_exact}, not its linearisation.} The
first-order form $\epsplat \approx |\overline{S}\,\Delta\ln u|$ is adequate for
a small step and wrong for a large one: at the four-interval step used below it
returns $59.4\,\%$ where the exact expression gives $81.1\,\%$, understating by
more than a quarter (\texttt{verify\_ridge\_mechanism.py}). The exponential
costs nothing.

One caution about what follows. Because Eq.~\eqref{eq:exact_decomposition} is an
identity, comparing its two sides tests nothing: a check that $\epsplat$ equals
$|\mathrm{e}^{\overline{S}\Delta\ln u}-1|$ can fail only through arithmetic
error. What the identity permits is an \emph{attribution} --- the two factors
are separately measurable, so one may ask which carries the structure the maps
show. That is a weaker kind of statement than a derivation, and it is the kind
made below.

\subsection{Across density: a sensitivity structure}
\label{sec:density_mechanism}

The plateau error has an interior maximum in density at
$\nel \approx 1.93\times10^{13}\,\si{\per\cubic\centi\metre}$, at every
temperature in the range this thesis defends.

Evaluating both factors of Eq.~\eqref{eq:exact_decomposition} at every grid
point settles which of them puts it there
(\texttt{verify\_ridge\_mechanism.py}, over $\Te = 2$--$9\,$eV and all eight
density columns). The column maximising $\epsplat$ is the column maximising
$|\overline{S}|$ in $32$ of $32$ temperature rows for a one-interval
temperature step and $31$ of $31$ for a four-interval step; the column
maximising $|\Delta\ln u|$ coincides with the error maximum in $0$ of $32$ and
$0$ of $31$.

Coincidence counts are weak evidence on their own, because a quantity that
barely varies still has its maximum somewhere. The amplitudes settle it. Across
the three decades of density $|\overline{S}|$ varies by a factor $7.11$ and
$|\Delta\ln u|$ by a factor $1.01$. \textbf{A quantity that changes by one per
cent cannot produce a ridge in which the error changes by a factor of seven},
wherever its own maximum falls. Table~\ref{tab:mechanism_profiles} makes the
point without statistics.

\begin{table}[htbp]
\centering
\caption{Density profiles of the plateau error and of the two factors of
Eq.~\eqref{eq:exact_decomposition}, each normalised to its own row maximum and
then taken as the median over the $32$ temperature rows with
$\Te = 2$--$9\,$eV, for a one-interval temperature step. From
\texttt{verify\_ridge\_mechanism.py}.}
\label{tab:mechanism_profiles}
\sisetup{table-format=1.3}
\begin{tabular}{lSSSSSSSS}
  \toprule
  $\nel$ (\si{\per\cubic\centi\metre}) & {$10^{12}$} & {$2.7{\cdot}10^{12}$} &
  {$7.2{\cdot}10^{12}$} & {$1.9{\cdot}10^{13}$} & {$5.2{\cdot}10^{13}$} &
  {$1.4{\cdot}10^{14}$} & {$3.7{\cdot}10^{14}$} & {$10^{15}$} \\
  \midrule
  $\epsplat$        & 0.437 & 0.683 & 0.937 & 1.000 & 0.807 & 0.512 & 0.280 & 0.138 \\
  $|\overline{S}|$  & 0.445 & 0.689 & 0.938 & 1.000 & 0.813 & 0.520 & 0.285 & 0.141 \\
  $|\Delta\ln u|$   & 0.997 & 0.995 & 0.993 & 0.991 & 0.990 & 0.992 & 0.996 & 1.000 \\
  \bottomrule
\end{tabular}
\end{table}

The first two rows agree to $2.2\,\%$ at every column; the third is flat to one
per cent. The density structure of the plateau error is the density structure of
the observable's sensitivity to its reservoir.

That sensitivity has an interior maximum for the reason
Section~\ref{sec:ground_fed_fraction} gives. Both $\ffrac{3}$ and $\ffrac{4}$
rise from zero to unity as the reservoir grows, but the two shells switch from
recombination-fed to ground-fed at different reservoir strengths, because the
ground-fed channel weakens with principal quantum number faster than the
recombination-fed one. Their difference therefore vanishes in both supply limits
and peaks between them. The density ridge is where the reservoir sits in that
mixed-supply region.

\subsection{Alternatives that would produce the same appearance}
\label{sec:ridge_alternatives}

An interior maximum on a coarse grid can be manufactured in several ways, and
the attribution above does not by itself exclude them.

\emph{Grid resolution.} The ridge occupies one density column and the grid is
spaced by a factor $2.683$. Its location is therefore known only to that factor,
and this work cannot say whether the true maximum lies at
$1.93\times10^{13}\,\si{\per\cubic\centi\metre}$ or anywhere within a factor
$2.7$ of it. What is established is that the maximum is interior --- not at
either end of the density range --- and that it does not move between
temperature rows. Refining the density grid requires rebuilding the rate matrix
and revalidating every number in Chapter~\ref{ch:theory} against the new one;
future work could determine the location more precisely.

\emph{Step-size convention.} A ridge whose position depended on the size of the
imposed step would be an artifact of that choice. It does not: the error maximum
falls on the same density column for a one-interval and a four-interval step in
$31$ of the $31$ temperature rows both can reach, and so does the maximum of
$|\overline{S}|$ (\texttt{verify\_ridge\_mechanism.py}). This tests steps of
fixed \emph{fractional} size; a step of fixed absolute size in electron-volts
spans a different number of grid intervals at each temperature and is a
different experiment, not tested here.

\emph{Truncation.} The model resolves $n\le8$ and bundles $n=9$--$15$. The
mechanism turns on the difference between the $n=3$ and $n=4$ supply fractions,
both fully resolved, but those fractions depend on the recombination flux
cascading from above, which the bundling approximates. This work does not test
the sensitivity of the ridge to $n_{\max}$; Chapter~\ref{ch:limits} lists it
among the outstanding checks.

\subsection{Across temperature: a displacement effect}
\label{sec:temperature_mechanism}

The temperature dependence has the opposite explanation, and this is where
treating the two axes alike would go wrong.

Along the ridge column the plateau error falls monotonically from $18.07\,\%$ at
$\Te = 2.02\,$eV to $4.55\,\%$ at $8.29\,$eV, a factor $3.72$. Decomposed
through Eq.~\eqref{eq:exact_decomposition}, $|\Delta\ln u|$ contributes a factor
$2.72$ of that fall and $|\overline{S}|$ only $1.29$; in logarithmic terms the
displacement carries $76\,\%$ of the decay and the sensitivity $19\,\%$. At the
four-interval step the split is $72\,\%$ and $13\,\%$
(\texttt{verify\_ridge\_mechanism.py}).

The observable therefore does not become much less sensitive as the plasma
heats. \textbf{The ground state simply moves less.} A step of fixed fractional
size displaces the reservoir progressively less at higher temperature, and the
error follows it down.

Why the displacement behaves this way is not established here. In particular
$|\Delta\ln u|$ is flat to one per cent across three decades of density while
varying by a factor $2.7$ across temperature, which implies that
$\partial\ln\ngs/\partial\ln\Te$ at collisional--radiative equilibrium is
controlled by temperature alone. \emph{[MECHANISM NOT ESTABLISHED: this is an
observation from \texttt{verify\_ridge\_mechanism.py}, not a derived result. It
has not been shown analytically and is a candidate for future work.]}

There is no interior maximum in temperature anywhere in the defended range. The
largest values, near $\Te\approx2\,$eV, are boundary values of the
optical-thickness restriction of Section~\ref{sec:optical_depth}, not a peak.

\subsection{Three explanations this work does not support}
\label{sec:withdrawn}

Three readings of the ridge are natural and do not survive the numbers. Each is
recorded because it is among the first a reader is likely to propose.

\textbf{It is not detachment.} The ridge sits at
$1.93\times10^{19}\,\si{\per\cubic\metre}$: a factor $5.2$ below the density
band that published edge modelling associates with a detached ITER divertor
\emph{[SOURCE REQUIRED: Guillemaut et al.\ 2011, Fus.\ Eng.\ Des.\ 86, 2954 ---
key not yet in the bibliography]}, and a factor $1.6$ below the modelled
upstream separatrix density \emph{[SOURCE REQUIRED: Stangeby et al.\ 2022,
Nucl.\ Fusion 62]}. It falls between the two, in a region neither source
characterises. The ridge is reproducible; its identification with a named
divertor operating regime is not supported by anything in this work.

\textbf{It is not a peak in temperature.} The error decreases monotonically with
temperature throughout the defended range, as
Section~\ref{sec:temperature_mechanism} shows.

\textbf{It is not simply where $\bdep^{\rm CRE}$ meets the sensitivity
maximum.} That reading accounts for the density structure but fails on the
temperature axis: at the four-interval step the ratio
$\bdep^{\rm CRE}/\bdep^{\rm peak}$ passes through unity at $\Te = 4.7\,$eV,
where the error is $33\,\%$ and falling monotonically
(\texttt{verify\_plateau\_gridmap.py}). The endpoint of the reservoir
trajectory is the wrong object: Eq.~\eqref{eq:Sbar_def} shows that what enters
is the mean sensitivity over the whole interval traversed.

\subsection{What this section establishes}

The plateau error factorises exactly into a sensitivity and a displacement,
Eq.~\eqref{eq:eps_plateau_exact}. Attributing the structure of the error maps to
one factor or the other gives different answers on the two axes: across density
the sensitivity carries the structure while the displacement is flat to one per
cent, and across temperature the displacement carries three quarters of the
decay. The ridge is a property of how the $n=3$ and $n=4$ shells divide their
supply between excitation from the ground state and recombination from the
continuum --- a property of the atomic system, needing no divertor geometry to
produce it --- while its amplitude at any density is set by how far a
temperature step moves the ground state.

That the local model generates a ridge without spatial structure does not imply
spatial structure is irrelevant to a real line-integrated measurement, which
Chapter~\ref{ch:limits} takes up.

The argument so far is entirely algebraic. Every quantity above is built from
the partial-equilibrium state $\Robs^{\rm PE}$, and that state was constructed
by freezing the ground state by hand rather than by integrating anything.
Whether the time-dependent system actually visits it is a separate question.


% ------------------------------------------------------------------
% Sections still to be written. The labels are declared here so that the
% cross-references in chapter5_C5D.tex resolve rather than printing "??".
% Replace each stub with the real section as it is drafted.
% ------------------------------------------------------------------
\section{Does the transient visit the plateau?}
\label{sec:bridge}

\emph{[TO BE WRITTEN --- C5-A writeup. Full 43-state integration; the
dynamic bridge; the closed form for the early-response slope $k$, which was
cut from Section~\ref{sec:mechanism} because it needs the transient
established first.]}

\section{How long the error persists}
\label{sec:persistence}

\emph{[TO BE WRITTEN. $\tau_{\rm eff}/\tauqss = 1.01$ and $1.07$.]}

\section{Optical thickness and the restricted range}
\label{sec:optical_depth}

\emph{[TO BE WRITTEN --- may instead live in Chapter~\ref{ch:theory};
Section~\ref{sec:mechanism} refers to it for the $\Te \ge 2$~eV restriction.]}
          % if the other chapters use \include
%   or
%     % ============================================================
% chapter5.tex -- Chapter 5 wrapper.
%
% chapter5_C5D.tex is a SECTION (\section{...}), not a chapter. It needs a
% \chapter{} above it or LaTeX has nowhere to attach it.
%
% Add to thesis_main.tex, in the document body, after chapter 4:
%
%     % ============================================================
% chapter5.tex -- Chapter 5 wrapper.
%
% chapter5_C5D.tex is a SECTION (\section{...}), not a chapter. It needs a
% \chapter{} above it or LaTeX has nowhere to attach it.
%
% Add to thesis_main.tex, in the document body, after chapter 4:
%
%     % ============================================================
% chapter5.tex -- Chapter 5 wrapper.
%
% chapter5_C5D.tex is a SECTION (\section{...}), not a chapter. It needs a
% \chapter{} above it or LaTeX has nowhere to attach it.
%
% Add to thesis_main.tex, in the document body, after chapter 4:
%
%     \include{chapter5}          % if the other chapters use \include
%   or
%     \input{chapter5}            % if they use \input
%
% Use whichever the neighbouring chapters use -- mixing them causes the page
% breaks and numbering to differ from the rest of the thesis.
% ============================================================

\chapter{What the Model Says}
\label{ch:results}

% ------------------------------------------------------------------
% PLACEHOLDER so the chapter is never empty while sections are drafted.
% Delete this block once sec:error_maps exists.
% ------------------------------------------------------------------
\section{The error maps}
\label{sec:error_maps}

\emph{[TO BE WRITTEN --- C5-F. Defines the operating grid, presents
$\epsstep$ and $\epsplat$ over it, and gives the headline numbers with their
scopes. Section~\ref{sec:mechanism} below is drafted and refers back to this
section for the maps it explains.]}

% ------------------------------------------------------------------
% C5-D, drafted 24 August 2026.
% ------------------------------------------------------------------
\input{chapter5_C5D}

% ------------------------------------------------------------------
% Sections still to be written. The labels are declared here so that the
% cross-references in chapter5_C5D.tex resolve rather than printing "??".
% Replace each stub with the real section as it is drafted.
% ------------------------------------------------------------------
\section{Does the transient visit the plateau?}
\label{sec:bridge}

\emph{[TO BE WRITTEN --- C5-A writeup. Full 43-state integration; the
dynamic bridge; the closed form for the early-response slope $k$, which was
cut from Section~\ref{sec:mechanism} because it needs the transient
established first.]}

\section{How long the error persists}
\label{sec:persistence}

\emph{[TO BE WRITTEN. $\tau_{\rm eff}/\tauqss = 1.01$ and $1.07$.]}

\section{Optical thickness and the restricted range}
\label{sec:optical_depth}

\emph{[TO BE WRITTEN --- may instead live in Chapter~\ref{ch:theory};
Section~\ref{sec:mechanism} refers to it for the $\Te \ge 2$~eV restriction.]}
          % if the other chapters use \include
%   or
%     % ============================================================
% chapter5.tex -- Chapter 5 wrapper.
%
% chapter5_C5D.tex is a SECTION (\section{...}), not a chapter. It needs a
% \chapter{} above it or LaTeX has nowhere to attach it.
%
% Add to thesis_main.tex, in the document body, after chapter 4:
%
%     \include{chapter5}          % if the other chapters use \include
%   or
%     \input{chapter5}            % if they use \input
%
% Use whichever the neighbouring chapters use -- mixing them causes the page
% breaks and numbering to differ from the rest of the thesis.
% ============================================================

\chapter{What the Model Says}
\label{ch:results}

% ------------------------------------------------------------------
% PLACEHOLDER so the chapter is never empty while sections are drafted.
% Delete this block once sec:error_maps exists.
% ------------------------------------------------------------------
\section{The error maps}
\label{sec:error_maps}

\emph{[TO BE WRITTEN --- C5-F. Defines the operating grid, presents
$\epsstep$ and $\epsplat$ over it, and gives the headline numbers with their
scopes. Section~\ref{sec:mechanism} below is drafted and refers back to this
section for the maps it explains.]}

% ------------------------------------------------------------------
% C5-D, drafted 24 August 2026.
% ------------------------------------------------------------------
\input{chapter5_C5D}

% ------------------------------------------------------------------
% Sections still to be written. The labels are declared here so that the
% cross-references in chapter5_C5D.tex resolve rather than printing "??".
% Replace each stub with the real section as it is drafted.
% ------------------------------------------------------------------
\section{Does the transient visit the plateau?}
\label{sec:bridge}

\emph{[TO BE WRITTEN --- C5-A writeup. Full 43-state integration; the
dynamic bridge; the closed form for the early-response slope $k$, which was
cut from Section~\ref{sec:mechanism} because it needs the transient
established first.]}

\section{How long the error persists}
\label{sec:persistence}

\emph{[TO BE WRITTEN. $\tau_{\rm eff}/\tauqss = 1.01$ and $1.07$.]}

\section{Optical thickness and the restricted range}
\label{sec:optical_depth}

\emph{[TO BE WRITTEN --- may instead live in Chapter~\ref{ch:theory};
Section~\ref{sec:mechanism} refers to it for the $\Te \ge 2$~eV restriction.]}
            % if they use \input
%
% Use whichever the neighbouring chapters use -- mixing them causes the page
% breaks and numbering to differ from the rest of the thesis.
% ============================================================

\chapter{What the Model Says}
\label{ch:results}

% ------------------------------------------------------------------
% PLACEHOLDER so the chapter is never empty while sections are drafted.
% Delete this block once sec:error_maps exists.
% ------------------------------------------------------------------
\section{The error maps}
\label{sec:error_maps}

\emph{[TO BE WRITTEN --- C5-F. Defines the operating grid, presents
$\epsstep$ and $\epsplat$ over it, and gives the headline numbers with their
scopes. Section~\ref{sec:mechanism} below is drafted and refers back to this
section for the maps it explains.]}

% ------------------------------------------------------------------
% C5-D, drafted 24 August 2026.
% ------------------------------------------------------------------
% ============================================================
% Chapter 5, Section: the mechanism.  C5-D, draft 2.
%
% Written under the thesis-writing skill, 24 August 2026.
% Source: CH5_EVIDENCE.md. L_grid SHA 2d92b58e..., S_grid SHA 7822f536...
% Draft 1 retained as chapter5_C5D_draft1_20260824.tex (failed six gates).
%
% MACROS beyond those chapter3.tex exercises:  \epsstep  \epsplat
%   (recorded in notation_and_definitions.md Sec. 7, never yet compiled --
%    CONFIRM they exist in thesis_main.tex before building.)
% ============================================================

%-----------------------------------------------------------------------------
\section{Why the error is largest where it is}
\label{sec:mechanism}
%-----------------------------------------------------------------------------

The maps of Section~\ref{sec:error_maps} show where the diagnostic error is
large. They do not say why, and a map without a mechanism is a lookup table
that cannot be extrapolated past the grid that produced it. This section asks
which physical quantity puts the ridge where it is.

The answer has two halves and they disagree. The structure along the density
axis and the structure along the temperature axis have different causes, and
reading them as one effect --- which the appearance of a single ridged surface
invites --- gets both wrong.

\subsection{Two error measures}
\label{sec:eps_definitions}

Section~\ref{sec:two_references} fixed the two reference states against which an
observable can be judged. Two further quantities are needed here, and both are
properties of a temperature \emph{step} rather than of a single operator, which
is why neither appears in Chapter~\ref{ch:theory}.

Let the plasma sit at the collisional--radiative equilibrium of
$\Lmat(\Te,\nel)$, and let the electron temperature step to $\Te+\Delta\Te$ at
$t=0$ with $\nel$ held fixed. Write $\Lmat^-$ and $\Lmat^+$ for the operators
before and after, in the notation of Section~\ref{sec:two_clocks}. The
\emph{step error}
\begin{equation}
  \epsstep = \left|\frac{\Rcre(\Te,\nel)}
                        {\Rcre(\Te+\Delta\Te,\nel)} - 1\right|
  \label{eq:eps_step_def}
\end{equation}

\noindent is the dimensionless distance between the two tabulated answers: what
the reading would change by if the plasma re-equilibrated instantly. The
\emph{plateau error}
\begin{equation}
  \epsplat = \left|\frac{\Robs^{\rm PE}}
                        {\Rcre(\Te+\Delta\Te,\nel)} - 1\right|,
  \qquad
  \Robs^{\rm PE} \equiv \Rqss\!\left(\Te+\Delta\Te,\nel;\,\bdep^-\right),
  \label{eq:eps_plateau_def}
\end{equation}

\noindent is the distance between the tabulated answer and the state the plasma
occupies once the excited manifold has settled but before the ground state has
moved. The superscript PE denotes \emph{partial equilibrium}: the post-step
excited response of Eq.~\eqref{eq:qss_ratio}, evaluated on the \emph{pre-step}
reservoir.

The consequence of keeping these separate is that they are not close. Across the
operating grid $\epsplat$ exceeds $\epsstep$ at every point tested, typically by
an order of magnitude (Section~\ref{sec:error_maps}). A diagnostician who
pictures the reading moving smoothly between the two tabulated values has not
merely misjudged the size of the excursion but has the wrong shape of transient.
$\epsplat$ is the quantity this chapter reports, and Section~\ref{sec:bridge}
establishes that the time-dependent system does pass through the state it
measures.

\subsection{An exact decomposition of the finite step}
\label{sec:exact_decomposition}

Chapter~\ref{ch:theory} worked in the ground-state departure coefficient
$\bdep = \ngs/(\Zsb\nel\nion)$, because that is the variable in which the
observable's response has the closed form of Eq.~\eqref{eq:R_of_b}. Here it is
better to use the reservoir itself,
\begin{equation}
  u \;\equiv\; \frac{\ngs}{\nion} \;=\; \bdep\,\Zsb(\Te)\,\nel ,
  \label{eq:u_def}
\end{equation}

\noindent a dimensionless ground-state population per unit ion density. At
fixed $(\Te,\nel)$ the two are proportional, so $\mathrm{d}\ln u =
\mathrm{d}\ln\bdep$ and the sensitivity of
Section~\ref{sec:ground_fed_fraction} carries over unchanged. The reason for
the change of variable is that a temperature step moves $\Zsb$: comparing a
pre-step $\bdep$ with a post-step one would mix a change in the reservoir with
a change in the Saha--Boltzmann normalisation, and $u$ has no normalisation to
move. That confusion is not hypothetical --- an earlier version of the analysis
reported in this section made exactly it, and the resulting contamination in
$\Delta\ln\bdep$ was comparable to the quantity being measured.

\textbf{Assumption.} Equation~\eqref{eq:u_def} treats $\nion$ as a fixed
reservoir, unchanged by the step, as throughout this thesis
(Section~\ref{sec:rate_equation}). This removes the freedom for the ionisation
balance to respond, and it is what makes $u$ numerically equal to the
ground-state population in the solved system. Chapter~\ref{ch:limits} returns
to what a mobile ion population would change.

With that, Eq.~\eqref{eq:R_of_b} makes the post-step observable a ratio of two
affine functions of $u$, so the logarithmic sensitivity of
Section~\ref{sec:ground_fed_fraction} is $\mathrm{d}\ln\Robs/\mathrm{d}\ln u =
\ffrac{3}-\ffrac{4}$, the difference of the two ground-fed fractions.
Integrating between the frozen and the equilibrated reservoir gives, with no
approximation,
\begin{equation}
  \ln\frac{\Robs^{\rm PE}}{\Rcre}
  \;=\; \int_{\ln u^+}^{\ln u^-}\!\!\big(\ffrac{3}-\ffrac{4}\big)\,\mathrm{d}\ln u
  \;=\; \overline{S}\;\Delta\ln u ,
  \label{eq:exact_decomposition}
\end{equation}

\noindent where $\Delta\ln u = \ln(u^-/u^+)$ is the reservoir displacement and
\begin{equation}
  \overline{S} \;\equiv\;
  \frac{1}{\Delta\ln u}\int_{\ln u^+}^{\ln u^-}\!\!
  \big(\ffrac{3}-\ffrac{4}\big)\,\mathrm{d}\ln u
  \label{eq:Sbar_def}
\end{equation}

\noindent is the mean sensitivity over the interval the reservoir traverses.
Both factors are dimensionless. Since $\ffrac{3},\ffrac{4}\in[0,1]$ by
Section~\ref{sec:ground_fed_fraction}, and $\overline{S}$ is an average of their
difference, $|\overline{S}|$ is bounded by the peak sensitivity of
Eq.~\eqref{eq:peak_sensitivity} and so is strictly less than unity everywhere.
The plateau error follows:
\begin{equation}
  \boxed{\ \epsplat = \left|\,\mathrm{e}^{\overline{S}\,\Delta\ln u} - 1\,\right|\ }
  \label{eq:eps_plateau_exact}
\end{equation}

The consequence worth extracting is that a single scalar error has become a
product of two quantities that can be measured independently: how strongly the
observable responds to its reservoir, and how far the reservoir moved. Nothing
in Eq.~\eqref{eq:exact_decomposition} says which of them varies across the
operating map, and that is what the rest of this section determines.

\textbf{Use Eq.~\eqref{eq:eps_plateau_exact}, not its linearisation.} The
first-order form $\epsplat \approx |\overline{S}\,\Delta\ln u|$ is adequate for
a small step and wrong for a large one: at the four-interval step used below it
returns $59.4\,\%$ where the exact expression gives $81.1\,\%$, understating by
more than a quarter (\texttt{verify\_ridge\_mechanism.py}). The exponential
costs nothing.

One caution about what follows. Because Eq.~\eqref{eq:exact_decomposition} is an
identity, comparing its two sides tests nothing: a check that $\epsplat$ equals
$|\mathrm{e}^{\overline{S}\Delta\ln u}-1|$ can fail only through arithmetic
error. What the identity permits is an \emph{attribution} --- the two factors
are separately measurable, so one may ask which carries the structure the maps
show. That is a weaker kind of statement than a derivation, and it is the kind
made below.

\subsection{Across density: a sensitivity structure}
\label{sec:density_mechanism}

The plateau error has an interior maximum in density at
$\nel \approx 1.93\times10^{13}\,\si{\per\cubic\centi\metre}$, at every
temperature in the range this thesis defends.

Evaluating both factors of Eq.~\eqref{eq:exact_decomposition} at every grid
point settles which of them puts it there
(\texttt{verify\_ridge\_mechanism.py}, over $\Te = 2$--$9\,$eV and all eight
density columns). The column maximising $\epsplat$ is the column maximising
$|\overline{S}|$ in $32$ of $32$ temperature rows for a one-interval
temperature step and $31$ of $31$ for a four-interval step; the column
maximising $|\Delta\ln u|$ coincides with the error maximum in $0$ of $32$ and
$0$ of $31$.

Coincidence counts are weak evidence on their own, because a quantity that
barely varies still has its maximum somewhere. The amplitudes settle it. Across
the three decades of density $|\overline{S}|$ varies by a factor $7.11$ and
$|\Delta\ln u|$ by a factor $1.01$. \textbf{A quantity that changes by one per
cent cannot produce a ridge in which the error changes by a factor of seven},
wherever its own maximum falls. Table~\ref{tab:mechanism_profiles} makes the
point without statistics.

\begin{table}[htbp]
\centering
\caption{Density profiles of the plateau error and of the two factors of
Eq.~\eqref{eq:exact_decomposition}, each normalised to its own row maximum and
then taken as the median over the $32$ temperature rows with
$\Te = 2$--$9\,$eV, for a one-interval temperature step. From
\texttt{verify\_ridge\_mechanism.py}.}
\label{tab:mechanism_profiles}
\sisetup{table-format=1.3}
\begin{tabular}{lSSSSSSSS}
  \toprule
  $\nel$ (\si{\per\cubic\centi\metre}) & {$10^{12}$} & {$2.7{\cdot}10^{12}$} &
  {$7.2{\cdot}10^{12}$} & {$1.9{\cdot}10^{13}$} & {$5.2{\cdot}10^{13}$} &
  {$1.4{\cdot}10^{14}$} & {$3.7{\cdot}10^{14}$} & {$10^{15}$} \\
  \midrule
  $\epsplat$        & 0.437 & 0.683 & 0.937 & 1.000 & 0.807 & 0.512 & 0.280 & 0.138 \\
  $|\overline{S}|$  & 0.445 & 0.689 & 0.938 & 1.000 & 0.813 & 0.520 & 0.285 & 0.141 \\
  $|\Delta\ln u|$   & 0.997 & 0.995 & 0.993 & 0.991 & 0.990 & 0.992 & 0.996 & 1.000 \\
  \bottomrule
\end{tabular}
\end{table}

The first two rows agree to $2.2\,\%$ at every column; the third is flat to one
per cent. The density structure of the plateau error is the density structure of
the observable's sensitivity to its reservoir.

That sensitivity has an interior maximum for the reason
Section~\ref{sec:ground_fed_fraction} gives. Both $\ffrac{3}$ and $\ffrac{4}$
rise from zero to unity as the reservoir grows, but the two shells switch from
recombination-fed to ground-fed at different reservoir strengths, because the
ground-fed channel weakens with principal quantum number faster than the
recombination-fed one. Their difference therefore vanishes in both supply limits
and peaks between them. The density ridge is where the reservoir sits in that
mixed-supply region.

\subsection{Alternatives that would produce the same appearance}
\label{sec:ridge_alternatives}

An interior maximum on a coarse grid can be manufactured in several ways, and
the attribution above does not by itself exclude them.

\emph{Grid resolution.} The ridge occupies one density column and the grid is
spaced by a factor $2.683$. Its location is therefore known only to that factor,
and this work cannot say whether the true maximum lies at
$1.93\times10^{13}\,\si{\per\cubic\centi\metre}$ or anywhere within a factor
$2.7$ of it. What is established is that the maximum is interior --- not at
either end of the density range --- and that it does not move between
temperature rows. Refining the density grid requires rebuilding the rate matrix
and revalidating every number in Chapter~\ref{ch:theory} against the new one;
future work could determine the location more precisely.

\emph{Step-size convention.} A ridge whose position depended on the size of the
imposed step would be an artifact of that choice. It does not: the error maximum
falls on the same density column for a one-interval and a four-interval step in
$31$ of the $31$ temperature rows both can reach, and so does the maximum of
$|\overline{S}|$ (\texttt{verify\_ridge\_mechanism.py}). This tests steps of
fixed \emph{fractional} size; a step of fixed absolute size in electron-volts
spans a different number of grid intervals at each temperature and is a
different experiment, not tested here.

\emph{Truncation.} The model resolves $n\le8$ and bundles $n=9$--$15$. The
mechanism turns on the difference between the $n=3$ and $n=4$ supply fractions,
both fully resolved, but those fractions depend on the recombination flux
cascading from above, which the bundling approximates. This work does not test
the sensitivity of the ridge to $n_{\max}$; Chapter~\ref{ch:limits} lists it
among the outstanding checks.

\subsection{Across temperature: a displacement effect}
\label{sec:temperature_mechanism}

The temperature dependence has the opposite explanation, and this is where
treating the two axes alike would go wrong.

Along the ridge column the plateau error falls monotonically from $18.07\,\%$ at
$\Te = 2.02\,$eV to $4.55\,\%$ at $8.29\,$eV, a factor $3.72$. Decomposed
through Eq.~\eqref{eq:exact_decomposition}, $|\Delta\ln u|$ contributes a factor
$2.72$ of that fall and $|\overline{S}|$ only $1.29$; in logarithmic terms the
displacement carries $76\,\%$ of the decay and the sensitivity $19\,\%$. At the
four-interval step the split is $72\,\%$ and $13\,\%$
(\texttt{verify\_ridge\_mechanism.py}).

The observable therefore does not become much less sensitive as the plasma
heats. \textbf{The ground state simply moves less.} A step of fixed fractional
size displaces the reservoir progressively less at higher temperature, and the
error follows it down.

Why the displacement behaves this way is not established here. In particular
$|\Delta\ln u|$ is flat to one per cent across three decades of density while
varying by a factor $2.7$ across temperature, which implies that
$\partial\ln\ngs/\partial\ln\Te$ at collisional--radiative equilibrium is
controlled by temperature alone. \emph{[MECHANISM NOT ESTABLISHED: this is an
observation from \texttt{verify\_ridge\_mechanism.py}, not a derived result. It
has not been shown analytically and is a candidate for future work.]}

There is no interior maximum in temperature anywhere in the defended range. The
largest values, near $\Te\approx2\,$eV, are boundary values of the
optical-thickness restriction of Section~\ref{sec:optical_depth}, not a peak.

\subsection{Three explanations this work does not support}
\label{sec:withdrawn}

Three readings of the ridge are natural and do not survive the numbers. Each is
recorded because it is among the first a reader is likely to propose.

\textbf{It is not detachment.} The ridge sits at
$1.93\times10^{19}\,\si{\per\cubic\metre}$: a factor $5.2$ below the density
band that published edge modelling associates with a detached ITER divertor
\emph{[SOURCE REQUIRED: Guillemaut et al.\ 2011, Fus.\ Eng.\ Des.\ 86, 2954 ---
key not yet in the bibliography]}, and a factor $1.6$ below the modelled
upstream separatrix density \emph{[SOURCE REQUIRED: Stangeby et al.\ 2022,
Nucl.\ Fusion 62]}. It falls between the two, in a region neither source
characterises. The ridge is reproducible; its identification with a named
divertor operating regime is not supported by anything in this work.

\textbf{It is not a peak in temperature.} The error decreases monotonically with
temperature throughout the defended range, as
Section~\ref{sec:temperature_mechanism} shows.

\textbf{It is not simply where $\bdep^{\rm CRE}$ meets the sensitivity
maximum.} That reading accounts for the density structure but fails on the
temperature axis: at the four-interval step the ratio
$\bdep^{\rm CRE}/\bdep^{\rm peak}$ passes through unity at $\Te = 4.7\,$eV,
where the error is $33\,\%$ and falling monotonically
(\texttt{verify\_plateau\_gridmap.py}). The endpoint of the reservoir
trajectory is the wrong object: Eq.~\eqref{eq:Sbar_def} shows that what enters
is the mean sensitivity over the whole interval traversed.

\subsection{What this section establishes}

The plateau error factorises exactly into a sensitivity and a displacement,
Eq.~\eqref{eq:eps_plateau_exact}. Attributing the structure of the error maps to
one factor or the other gives different answers on the two axes: across density
the sensitivity carries the structure while the displacement is flat to one per
cent, and across temperature the displacement carries three quarters of the
decay. The ridge is a property of how the $n=3$ and $n=4$ shells divide their
supply between excitation from the ground state and recombination from the
continuum --- a property of the atomic system, needing no divertor geometry to
produce it --- while its amplitude at any density is set by how far a
temperature step moves the ground state.

That the local model generates a ridge without spatial structure does not imply
spatial structure is irrelevant to a real line-integrated measurement, which
Chapter~\ref{ch:limits} takes up.

The argument so far is entirely algebraic. Every quantity above is built from
the partial-equilibrium state $\Robs^{\rm PE}$, and that state was constructed
by freezing the ground state by hand rather than by integrating anything.
Whether the time-dependent system actually visits it is a separate question.


% ------------------------------------------------------------------
% Sections still to be written. The labels are declared here so that the
% cross-references in chapter5_C5D.tex resolve rather than printing "??".
% Replace each stub with the real section as it is drafted.
% ------------------------------------------------------------------
\section{Does the transient visit the plateau?}
\label{sec:bridge}

\emph{[TO BE WRITTEN --- C5-A writeup. Full 43-state integration; the
dynamic bridge; the closed form for the early-response slope $k$, which was
cut from Section~\ref{sec:mechanism} because it needs the transient
established first.]}

\section{How long the error persists}
\label{sec:persistence}

\emph{[TO BE WRITTEN. $\tau_{\rm eff}/\tauqss = 1.01$ and $1.07$.]}

\section{Optical thickness and the restricted range}
\label{sec:optical_depth}

\emph{[TO BE WRITTEN --- may instead live in Chapter~\ref{ch:theory};
Section~\ref{sec:mechanism} refers to it for the $\Te \ge 2$~eV restriction.]}
          % if the other chapters use \include
%   or
%     % ============================================================
% chapter5.tex -- Chapter 5 wrapper.
%
% chapter5_C5D.tex is a SECTION (\section{...}), not a chapter. It needs a
% \chapter{} above it or LaTeX has nowhere to attach it.
%
% Add to thesis_main.tex, in the document body, after chapter 4:
%
%     % ============================================================
% chapter5.tex -- Chapter 5 wrapper.
%
% chapter5_C5D.tex is a SECTION (\section{...}), not a chapter. It needs a
% \chapter{} above it or LaTeX has nowhere to attach it.
%
% Add to thesis_main.tex, in the document body, after chapter 4:
%
%     \include{chapter5}          % if the other chapters use \include
%   or
%     \input{chapter5}            % if they use \input
%
% Use whichever the neighbouring chapters use -- mixing them causes the page
% breaks and numbering to differ from the rest of the thesis.
% ============================================================

\chapter{What the Model Says}
\label{ch:results}

% ------------------------------------------------------------------
% PLACEHOLDER so the chapter is never empty while sections are drafted.
% Delete this block once sec:error_maps exists.
% ------------------------------------------------------------------
\section{The error maps}
\label{sec:error_maps}

\emph{[TO BE WRITTEN --- C5-F. Defines the operating grid, presents
$\epsstep$ and $\epsplat$ over it, and gives the headline numbers with their
scopes. Section~\ref{sec:mechanism} below is drafted and refers back to this
section for the maps it explains.]}

% ------------------------------------------------------------------
% C5-D, drafted 24 August 2026.
% ------------------------------------------------------------------
\input{chapter5_C5D}

% ------------------------------------------------------------------
% Sections still to be written. The labels are declared here so that the
% cross-references in chapter5_C5D.tex resolve rather than printing "??".
% Replace each stub with the real section as it is drafted.
% ------------------------------------------------------------------
\section{Does the transient visit the plateau?}
\label{sec:bridge}

\emph{[TO BE WRITTEN --- C5-A writeup. Full 43-state integration; the
dynamic bridge; the closed form for the early-response slope $k$, which was
cut from Section~\ref{sec:mechanism} because it needs the transient
established first.]}

\section{How long the error persists}
\label{sec:persistence}

\emph{[TO BE WRITTEN. $\tau_{\rm eff}/\tauqss = 1.01$ and $1.07$.]}

\section{Optical thickness and the restricted range}
\label{sec:optical_depth}

\emph{[TO BE WRITTEN --- may instead live in Chapter~\ref{ch:theory};
Section~\ref{sec:mechanism} refers to it for the $\Te \ge 2$~eV restriction.]}
          % if the other chapters use \include
%   or
%     % ============================================================
% chapter5.tex -- Chapter 5 wrapper.
%
% chapter5_C5D.tex is a SECTION (\section{...}), not a chapter. It needs a
% \chapter{} above it or LaTeX has nowhere to attach it.
%
% Add to thesis_main.tex, in the document body, after chapter 4:
%
%     \include{chapter5}          % if the other chapters use \include
%   or
%     \input{chapter5}            % if they use \input
%
% Use whichever the neighbouring chapters use -- mixing them causes the page
% breaks and numbering to differ from the rest of the thesis.
% ============================================================

\chapter{What the Model Says}
\label{ch:results}

% ------------------------------------------------------------------
% PLACEHOLDER so the chapter is never empty while sections are drafted.
% Delete this block once sec:error_maps exists.
% ------------------------------------------------------------------
\section{The error maps}
\label{sec:error_maps}

\emph{[TO BE WRITTEN --- C5-F. Defines the operating grid, presents
$\epsstep$ and $\epsplat$ over it, and gives the headline numbers with their
scopes. Section~\ref{sec:mechanism} below is drafted and refers back to this
section for the maps it explains.]}

% ------------------------------------------------------------------
% C5-D, drafted 24 August 2026.
% ------------------------------------------------------------------
\input{chapter5_C5D}

% ------------------------------------------------------------------
% Sections still to be written. The labels are declared here so that the
% cross-references in chapter5_C5D.tex resolve rather than printing "??".
% Replace each stub with the real section as it is drafted.
% ------------------------------------------------------------------
\section{Does the transient visit the plateau?}
\label{sec:bridge}

\emph{[TO BE WRITTEN --- C5-A writeup. Full 43-state integration; the
dynamic bridge; the closed form for the early-response slope $k$, which was
cut from Section~\ref{sec:mechanism} because it needs the transient
established first.]}

\section{How long the error persists}
\label{sec:persistence}

\emph{[TO BE WRITTEN. $\tau_{\rm eff}/\tauqss = 1.01$ and $1.07$.]}

\section{Optical thickness and the restricted range}
\label{sec:optical_depth}

\emph{[TO BE WRITTEN --- may instead live in Chapter~\ref{ch:theory};
Section~\ref{sec:mechanism} refers to it for the $\Te \ge 2$~eV restriction.]}
            % if they use \input
%
% Use whichever the neighbouring chapters use -- mixing them causes the page
% breaks and numbering to differ from the rest of the thesis.
% ============================================================

\chapter{What the Model Says}
\label{ch:results}

% ------------------------------------------------------------------
% PLACEHOLDER so the chapter is never empty while sections are drafted.
% Delete this block once sec:error_maps exists.
% ------------------------------------------------------------------
\section{The error maps}
\label{sec:error_maps}

\emph{[TO BE WRITTEN --- C5-F. Defines the operating grid, presents
$\epsstep$ and $\epsplat$ over it, and gives the headline numbers with their
scopes. Section~\ref{sec:mechanism} below is drafted and refers back to this
section for the maps it explains.]}

% ------------------------------------------------------------------
% C5-D, drafted 24 August 2026.
% ------------------------------------------------------------------
% ============================================================
% Chapter 5, Section: the mechanism.  C5-D, draft 2.
%
% Written under the thesis-writing skill, 24 August 2026.
% Source: CH5_EVIDENCE.md. L_grid SHA 2d92b58e..., S_grid SHA 7822f536...
% Draft 1 retained as chapter5_C5D_draft1_20260824.tex (failed six gates).
%
% MACROS beyond those chapter3.tex exercises:  \epsstep  \epsplat
%   (recorded in notation_and_definitions.md Sec. 7, never yet compiled --
%    CONFIRM they exist in thesis_main.tex before building.)
% ============================================================

%-----------------------------------------------------------------------------
\section{Why the error is largest where it is}
\label{sec:mechanism}
%-----------------------------------------------------------------------------

The maps of Section~\ref{sec:error_maps} show where the diagnostic error is
large. They do not say why, and a map without a mechanism is a lookup table
that cannot be extrapolated past the grid that produced it. This section asks
which physical quantity puts the ridge where it is.

The answer has two halves and they disagree. The structure along the density
axis and the structure along the temperature axis have different causes, and
reading them as one effect --- which the appearance of a single ridged surface
invites --- gets both wrong.

\subsection{Two error measures}
\label{sec:eps_definitions}

Section~\ref{sec:two_references} fixed the two reference states against which an
observable can be judged. Two further quantities are needed here, and both are
properties of a temperature \emph{step} rather than of a single operator, which
is why neither appears in Chapter~\ref{ch:theory}.

Let the plasma sit at the collisional--radiative equilibrium of
$\Lmat(\Te,\nel)$, and let the electron temperature step to $\Te+\Delta\Te$ at
$t=0$ with $\nel$ held fixed. Write $\Lmat^-$ and $\Lmat^+$ for the operators
before and after, in the notation of Section~\ref{sec:two_clocks}. The
\emph{step error}
\begin{equation}
  \epsstep = \left|\frac{\Rcre(\Te,\nel)}
                        {\Rcre(\Te+\Delta\Te,\nel)} - 1\right|
  \label{eq:eps_step_def}
\end{equation}

\noindent is the dimensionless distance between the two tabulated answers: what
the reading would change by if the plasma re-equilibrated instantly. The
\emph{plateau error}
\begin{equation}
  \epsplat = \left|\frac{\Robs^{\rm PE}}
                        {\Rcre(\Te+\Delta\Te,\nel)} - 1\right|,
  \qquad
  \Robs^{\rm PE} \equiv \Rqss\!\left(\Te+\Delta\Te,\nel;\,\bdep^-\right),
  \label{eq:eps_plateau_def}
\end{equation}

\noindent is the distance between the tabulated answer and the state the plasma
occupies once the excited manifold has settled but before the ground state has
moved. The superscript PE denotes \emph{partial equilibrium}: the post-step
excited response of Eq.~\eqref{eq:qss_ratio}, evaluated on the \emph{pre-step}
reservoir.

The consequence of keeping these separate is that they are not close. Across the
operating grid $\epsplat$ exceeds $\epsstep$ at every point tested, typically by
an order of magnitude (Section~\ref{sec:error_maps}). A diagnostician who
pictures the reading moving smoothly between the two tabulated values has not
merely misjudged the size of the excursion but has the wrong shape of transient.
$\epsplat$ is the quantity this chapter reports, and Section~\ref{sec:bridge}
establishes that the time-dependent system does pass through the state it
measures.

\subsection{An exact decomposition of the finite step}
\label{sec:exact_decomposition}

Chapter~\ref{ch:theory} worked in the ground-state departure coefficient
$\bdep = \ngs/(\Zsb\nel\nion)$, because that is the variable in which the
observable's response has the closed form of Eq.~\eqref{eq:R_of_b}. Here it is
better to use the reservoir itself,
\begin{equation}
  u \;\equiv\; \frac{\ngs}{\nion} \;=\; \bdep\,\Zsb(\Te)\,\nel ,
  \label{eq:u_def}
\end{equation}

\noindent a dimensionless ground-state population per unit ion density. At
fixed $(\Te,\nel)$ the two are proportional, so $\mathrm{d}\ln u =
\mathrm{d}\ln\bdep$ and the sensitivity of
Section~\ref{sec:ground_fed_fraction} carries over unchanged. The reason for
the change of variable is that a temperature step moves $\Zsb$: comparing a
pre-step $\bdep$ with a post-step one would mix a change in the reservoir with
a change in the Saha--Boltzmann normalisation, and $u$ has no normalisation to
move. That confusion is not hypothetical --- an earlier version of the analysis
reported in this section made exactly it, and the resulting contamination in
$\Delta\ln\bdep$ was comparable to the quantity being measured.

\textbf{Assumption.} Equation~\eqref{eq:u_def} treats $\nion$ as a fixed
reservoir, unchanged by the step, as throughout this thesis
(Section~\ref{sec:rate_equation}). This removes the freedom for the ionisation
balance to respond, and it is what makes $u$ numerically equal to the
ground-state population in the solved system. Chapter~\ref{ch:limits} returns
to what a mobile ion population would change.

With that, Eq.~\eqref{eq:R_of_b} makes the post-step observable a ratio of two
affine functions of $u$, so the logarithmic sensitivity of
Section~\ref{sec:ground_fed_fraction} is $\mathrm{d}\ln\Robs/\mathrm{d}\ln u =
\ffrac{3}-\ffrac{4}$, the difference of the two ground-fed fractions.
Integrating between the frozen and the equilibrated reservoir gives, with no
approximation,
\begin{equation}
  \ln\frac{\Robs^{\rm PE}}{\Rcre}
  \;=\; \int_{\ln u^+}^{\ln u^-}\!\!\big(\ffrac{3}-\ffrac{4}\big)\,\mathrm{d}\ln u
  \;=\; \overline{S}\;\Delta\ln u ,
  \label{eq:exact_decomposition}
\end{equation}

\noindent where $\Delta\ln u = \ln(u^-/u^+)$ is the reservoir displacement and
\begin{equation}
  \overline{S} \;\equiv\;
  \frac{1}{\Delta\ln u}\int_{\ln u^+}^{\ln u^-}\!\!
  \big(\ffrac{3}-\ffrac{4}\big)\,\mathrm{d}\ln u
  \label{eq:Sbar_def}
\end{equation}

\noindent is the mean sensitivity over the interval the reservoir traverses.
Both factors are dimensionless. Since $\ffrac{3},\ffrac{4}\in[0,1]$ by
Section~\ref{sec:ground_fed_fraction}, and $\overline{S}$ is an average of their
difference, $|\overline{S}|$ is bounded by the peak sensitivity of
Eq.~\eqref{eq:peak_sensitivity} and so is strictly less than unity everywhere.
The plateau error follows:
\begin{equation}
  \boxed{\ \epsplat = \left|\,\mathrm{e}^{\overline{S}\,\Delta\ln u} - 1\,\right|\ }
  \label{eq:eps_plateau_exact}
\end{equation}

The consequence worth extracting is that a single scalar error has become a
product of two quantities that can be measured independently: how strongly the
observable responds to its reservoir, and how far the reservoir moved. Nothing
in Eq.~\eqref{eq:exact_decomposition} says which of them varies across the
operating map, and that is what the rest of this section determines.

\textbf{Use Eq.~\eqref{eq:eps_plateau_exact}, not its linearisation.} The
first-order form $\epsplat \approx |\overline{S}\,\Delta\ln u|$ is adequate for
a small step and wrong for a large one: at the four-interval step used below it
returns $59.4\,\%$ where the exact expression gives $81.1\,\%$, understating by
more than a quarter (\texttt{verify\_ridge\_mechanism.py}). The exponential
costs nothing.

One caution about what follows. Because Eq.~\eqref{eq:exact_decomposition} is an
identity, comparing its two sides tests nothing: a check that $\epsplat$ equals
$|\mathrm{e}^{\overline{S}\Delta\ln u}-1|$ can fail only through arithmetic
error. What the identity permits is an \emph{attribution} --- the two factors
are separately measurable, so one may ask which carries the structure the maps
show. That is a weaker kind of statement than a derivation, and it is the kind
made below.

\subsection{Across density: a sensitivity structure}
\label{sec:density_mechanism}

The plateau error has an interior maximum in density at
$\nel \approx 1.93\times10^{13}\,\si{\per\cubic\centi\metre}$, at every
temperature in the range this thesis defends.

Evaluating both factors of Eq.~\eqref{eq:exact_decomposition} at every grid
point settles which of them puts it there
(\texttt{verify\_ridge\_mechanism.py}, over $\Te = 2$--$9\,$eV and all eight
density columns). The column maximising $\epsplat$ is the column maximising
$|\overline{S}|$ in $32$ of $32$ temperature rows for a one-interval
temperature step and $31$ of $31$ for a four-interval step; the column
maximising $|\Delta\ln u|$ coincides with the error maximum in $0$ of $32$ and
$0$ of $31$.

Coincidence counts are weak evidence on their own, because a quantity that
barely varies still has its maximum somewhere. The amplitudes settle it. Across
the three decades of density $|\overline{S}|$ varies by a factor $7.11$ and
$|\Delta\ln u|$ by a factor $1.01$. \textbf{A quantity that changes by one per
cent cannot produce a ridge in which the error changes by a factor of seven},
wherever its own maximum falls. Table~\ref{tab:mechanism_profiles} makes the
point without statistics.

\begin{table}[htbp]
\centering
\caption{Density profiles of the plateau error and of the two factors of
Eq.~\eqref{eq:exact_decomposition}, each normalised to its own row maximum and
then taken as the median over the $32$ temperature rows with
$\Te = 2$--$9\,$eV, for a one-interval temperature step. From
\texttt{verify\_ridge\_mechanism.py}.}
\label{tab:mechanism_profiles}
\sisetup{table-format=1.3}
\begin{tabular}{lSSSSSSSS}
  \toprule
  $\nel$ (\si{\per\cubic\centi\metre}) & {$10^{12}$} & {$2.7{\cdot}10^{12}$} &
  {$7.2{\cdot}10^{12}$} & {$1.9{\cdot}10^{13}$} & {$5.2{\cdot}10^{13}$} &
  {$1.4{\cdot}10^{14}$} & {$3.7{\cdot}10^{14}$} & {$10^{15}$} \\
  \midrule
  $\epsplat$        & 0.437 & 0.683 & 0.937 & 1.000 & 0.807 & 0.512 & 0.280 & 0.138 \\
  $|\overline{S}|$  & 0.445 & 0.689 & 0.938 & 1.000 & 0.813 & 0.520 & 0.285 & 0.141 \\
  $|\Delta\ln u|$   & 0.997 & 0.995 & 0.993 & 0.991 & 0.990 & 0.992 & 0.996 & 1.000 \\
  \bottomrule
\end{tabular}
\end{table}

The first two rows agree to $2.2\,\%$ at every column; the third is flat to one
per cent. The density structure of the plateau error is the density structure of
the observable's sensitivity to its reservoir.

That sensitivity has an interior maximum for the reason
Section~\ref{sec:ground_fed_fraction} gives. Both $\ffrac{3}$ and $\ffrac{4}$
rise from zero to unity as the reservoir grows, but the two shells switch from
recombination-fed to ground-fed at different reservoir strengths, because the
ground-fed channel weakens with principal quantum number faster than the
recombination-fed one. Their difference therefore vanishes in both supply limits
and peaks between them. The density ridge is where the reservoir sits in that
mixed-supply region.

\subsection{Alternatives that would produce the same appearance}
\label{sec:ridge_alternatives}

An interior maximum on a coarse grid can be manufactured in several ways, and
the attribution above does not by itself exclude them.

\emph{Grid resolution.} The ridge occupies one density column and the grid is
spaced by a factor $2.683$. Its location is therefore known only to that factor,
and this work cannot say whether the true maximum lies at
$1.93\times10^{13}\,\si{\per\cubic\centi\metre}$ or anywhere within a factor
$2.7$ of it. What is established is that the maximum is interior --- not at
either end of the density range --- and that it does not move between
temperature rows. Refining the density grid requires rebuilding the rate matrix
and revalidating every number in Chapter~\ref{ch:theory} against the new one;
future work could determine the location more precisely.

\emph{Step-size convention.} A ridge whose position depended on the size of the
imposed step would be an artifact of that choice. It does not: the error maximum
falls on the same density column for a one-interval and a four-interval step in
$31$ of the $31$ temperature rows both can reach, and so does the maximum of
$|\overline{S}|$ (\texttt{verify\_ridge\_mechanism.py}). This tests steps of
fixed \emph{fractional} size; a step of fixed absolute size in electron-volts
spans a different number of grid intervals at each temperature and is a
different experiment, not tested here.

\emph{Truncation.} The model resolves $n\le8$ and bundles $n=9$--$15$. The
mechanism turns on the difference between the $n=3$ and $n=4$ supply fractions,
both fully resolved, but those fractions depend on the recombination flux
cascading from above, which the bundling approximates. This work does not test
the sensitivity of the ridge to $n_{\max}$; Chapter~\ref{ch:limits} lists it
among the outstanding checks.

\subsection{Across temperature: a displacement effect}
\label{sec:temperature_mechanism}

The temperature dependence has the opposite explanation, and this is where
treating the two axes alike would go wrong.

Along the ridge column the plateau error falls monotonically from $18.07\,\%$ at
$\Te = 2.02\,$eV to $4.55\,\%$ at $8.29\,$eV, a factor $3.72$. Decomposed
through Eq.~\eqref{eq:exact_decomposition}, $|\Delta\ln u|$ contributes a factor
$2.72$ of that fall and $|\overline{S}|$ only $1.29$; in logarithmic terms the
displacement carries $76\,\%$ of the decay and the sensitivity $19\,\%$. At the
four-interval step the split is $72\,\%$ and $13\,\%$
(\texttt{verify\_ridge\_mechanism.py}).

The observable therefore does not become much less sensitive as the plasma
heats. \textbf{The ground state simply moves less.} A step of fixed fractional
size displaces the reservoir progressively less at higher temperature, and the
error follows it down.

Why the displacement behaves this way is not established here. In particular
$|\Delta\ln u|$ is flat to one per cent across three decades of density while
varying by a factor $2.7$ across temperature, which implies that
$\partial\ln\ngs/\partial\ln\Te$ at collisional--radiative equilibrium is
controlled by temperature alone. \emph{[MECHANISM NOT ESTABLISHED: this is an
observation from \texttt{verify\_ridge\_mechanism.py}, not a derived result. It
has not been shown analytically and is a candidate for future work.]}

There is no interior maximum in temperature anywhere in the defended range. The
largest values, near $\Te\approx2\,$eV, are boundary values of the
optical-thickness restriction of Section~\ref{sec:optical_depth}, not a peak.

\subsection{Three explanations this work does not support}
\label{sec:withdrawn}

Three readings of the ridge are natural and do not survive the numbers. Each is
recorded because it is among the first a reader is likely to propose.

\textbf{It is not detachment.} The ridge sits at
$1.93\times10^{19}\,\si{\per\cubic\metre}$: a factor $5.2$ below the density
band that published edge modelling associates with a detached ITER divertor
\emph{[SOURCE REQUIRED: Guillemaut et al.\ 2011, Fus.\ Eng.\ Des.\ 86, 2954 ---
key not yet in the bibliography]}, and a factor $1.6$ below the modelled
upstream separatrix density \emph{[SOURCE REQUIRED: Stangeby et al.\ 2022,
Nucl.\ Fusion 62]}. It falls between the two, in a region neither source
characterises. The ridge is reproducible; its identification with a named
divertor operating regime is not supported by anything in this work.

\textbf{It is not a peak in temperature.} The error decreases monotonically with
temperature throughout the defended range, as
Section~\ref{sec:temperature_mechanism} shows.

\textbf{It is not simply where $\bdep^{\rm CRE}$ meets the sensitivity
maximum.} That reading accounts for the density structure but fails on the
temperature axis: at the four-interval step the ratio
$\bdep^{\rm CRE}/\bdep^{\rm peak}$ passes through unity at $\Te = 4.7\,$eV,
where the error is $33\,\%$ and falling monotonically
(\texttt{verify\_plateau\_gridmap.py}). The endpoint of the reservoir
trajectory is the wrong object: Eq.~\eqref{eq:Sbar_def} shows that what enters
is the mean sensitivity over the whole interval traversed.

\subsection{What this section establishes}

The plateau error factorises exactly into a sensitivity and a displacement,
Eq.~\eqref{eq:eps_plateau_exact}. Attributing the structure of the error maps to
one factor or the other gives different answers on the two axes: across density
the sensitivity carries the structure while the displacement is flat to one per
cent, and across temperature the displacement carries three quarters of the
decay. The ridge is a property of how the $n=3$ and $n=4$ shells divide their
supply between excitation from the ground state and recombination from the
continuum --- a property of the atomic system, needing no divertor geometry to
produce it --- while its amplitude at any density is set by how far a
temperature step moves the ground state.

That the local model generates a ridge without spatial structure does not imply
spatial structure is irrelevant to a real line-integrated measurement, which
Chapter~\ref{ch:limits} takes up.

The argument so far is entirely algebraic. Every quantity above is built from
the partial-equilibrium state $\Robs^{\rm PE}$, and that state was constructed
by freezing the ground state by hand rather than by integrating anything.
Whether the time-dependent system actually visits it is a separate question.


% ------------------------------------------------------------------
% Sections still to be written. The labels are declared here so that the
% cross-references in chapter5_C5D.tex resolve rather than printing "??".
% Replace each stub with the real section as it is drafted.
% ------------------------------------------------------------------
\section{Does the transient visit the plateau?}
\label{sec:bridge}

\emph{[TO BE WRITTEN --- C5-A writeup. Full 43-state integration; the
dynamic bridge; the closed form for the early-response slope $k$, which was
cut from Section~\ref{sec:mechanism} because it needs the transient
established first.]}

\section{How long the error persists}
\label{sec:persistence}

\emph{[TO BE WRITTEN. $\tau_{\rm eff}/\tauqss = 1.01$ and $1.07$.]}

\section{Optical thickness and the restricted range}
\label{sec:optical_depth}

\emph{[TO BE WRITTEN --- may instead live in Chapter~\ref{ch:theory};
Section~\ref{sec:mechanism} refers to it for the $\Te \ge 2$~eV restriction.]}
            % if they use \input
%
% Use whichever the neighbouring chapters use -- mixing them causes the page
% breaks and numbering to differ from the rest of the thesis.
% ============================================================

\chapter{What the Model Says}
\label{ch:results}

% ------------------------------------------------------------------
% PLACEHOLDER so the chapter is never empty while sections are drafted.
% Delete this block once sec:error_maps exists.
% ------------------------------------------------------------------
\section{The error maps}
\label{sec:error_maps}

\emph{[TO BE WRITTEN --- C5-F. Defines the operating grid, presents
$\epsstep$ and $\epsplat$ over it, and gives the headline numbers with their
scopes. Section~\ref{sec:mechanism} below is drafted and refers back to this
section for the maps it explains.]}

% ------------------------------------------------------------------
% C5-D, drafted 24 August 2026.
% ------------------------------------------------------------------
% ============================================================
% Chapter 5, Section: the mechanism.  C5-D, draft 2.
%
% Written under the thesis-writing skill, 24 August 2026.
% Source: CH5_EVIDENCE.md. L_grid SHA 2d92b58e..., S_grid SHA 7822f536...
% Draft 1 retained as chapter5_C5D_draft1_20260824.tex (failed six gates).
%
% MACROS beyond those chapter3.tex exercises:  \epsstep  \epsplat
%   (recorded in notation_and_definitions.md Sec. 7, never yet compiled --
%    CONFIRM they exist in thesis_main.tex before building.)
% ============================================================

%-----------------------------------------------------------------------------
\section{Why the error is largest where it is}
\label{sec:mechanism}
%-----------------------------------------------------------------------------

The maps of Section~\ref{sec:error_maps} show where the diagnostic error is
large. They do not say why, and a map without a mechanism is a lookup table
that cannot be extrapolated past the grid that produced it. This section asks
which physical quantity puts the ridge where it is.

The answer has two halves and they disagree. The structure along the density
axis and the structure along the temperature axis have different causes, and
reading them as one effect --- which the appearance of a single ridged surface
invites --- gets both wrong.

\subsection{Two error measures}
\label{sec:eps_definitions}

Section~\ref{sec:two_references} fixed the two reference states against which an
observable can be judged. Two further quantities are needed here, and both are
properties of a temperature \emph{step} rather than of a single operator, which
is why neither appears in Chapter~\ref{ch:theory}.

Let the plasma sit at the collisional--radiative equilibrium of
$\Lmat(\Te,\nel)$, and let the electron temperature step to $\Te+\Delta\Te$ at
$t=0$ with $\nel$ held fixed. Write $\Lmat^-$ and $\Lmat^+$ for the operators
before and after, in the notation of Section~\ref{sec:two_clocks}. The
\emph{step error}
\begin{equation}
  \epsstep = \left|\frac{\Rcre(\Te,\nel)}
                        {\Rcre(\Te+\Delta\Te,\nel)} - 1\right|
  \label{eq:eps_step_def}
\end{equation}

\noindent is the dimensionless distance between the two tabulated answers: what
the reading would change by if the plasma re-equilibrated instantly. The
\emph{plateau error}
\begin{equation}
  \epsplat = \left|\frac{\Robs^{\rm PE}}
                        {\Rcre(\Te+\Delta\Te,\nel)} - 1\right|,
  \qquad
  \Robs^{\rm PE} \equiv \Rqss\!\left(\Te+\Delta\Te,\nel;\,\bdep^-\right),
  \label{eq:eps_plateau_def}
\end{equation}

\noindent is the distance between the tabulated answer and the state the plasma
occupies once the excited manifold has settled but before the ground state has
moved. The superscript PE denotes \emph{partial equilibrium}: the post-step
excited response of Eq.~\eqref{eq:qss_ratio}, evaluated on the \emph{pre-step}
reservoir.

The consequence of keeping these separate is that they are not close. Across the
operating grid $\epsplat$ exceeds $\epsstep$ at every point tested, typically by
an order of magnitude (Section~\ref{sec:error_maps}). A diagnostician who
pictures the reading moving smoothly between the two tabulated values has not
merely misjudged the size of the excursion but has the wrong shape of transient.
$\epsplat$ is the quantity this chapter reports, and Section~\ref{sec:bridge}
establishes that the time-dependent system does pass through the state it
measures.

\subsection{An exact decomposition of the finite step}
\label{sec:exact_decomposition}

Chapter~\ref{ch:theory} worked in the ground-state departure coefficient
$\bdep = \ngs/(\Zsb\nel\nion)$, because that is the variable in which the
observable's response has the closed form of Eq.~\eqref{eq:R_of_b}. Here it is
better to use the reservoir itself,
\begin{equation}
  u \;\equiv\; \frac{\ngs}{\nion} \;=\; \bdep\,\Zsb(\Te)\,\nel ,
  \label{eq:u_def}
\end{equation}

\noindent a dimensionless ground-state population per unit ion density. At
fixed $(\Te,\nel)$ the two are proportional, so $\mathrm{d}\ln u =
\mathrm{d}\ln\bdep$ and the sensitivity of
Section~\ref{sec:ground_fed_fraction} carries over unchanged. The reason for
the change of variable is that a temperature step moves $\Zsb$: comparing a
pre-step $\bdep$ with a post-step one would mix a change in the reservoir with
a change in the Saha--Boltzmann normalisation, and $u$ has no normalisation to
move. That confusion is not hypothetical --- an earlier version of the analysis
reported in this section made exactly it, and the resulting contamination in
$\Delta\ln\bdep$ was comparable to the quantity being measured.

\textbf{Assumption.} Equation~\eqref{eq:u_def} treats $\nion$ as a fixed
reservoir, unchanged by the step, as throughout this thesis
(Section~\ref{sec:rate_equation}). This removes the freedom for the ionisation
balance to respond, and it is what makes $u$ numerically equal to the
ground-state population in the solved system. Chapter~\ref{ch:limits} returns
to what a mobile ion population would change.

With that, Eq.~\eqref{eq:R_of_b} makes the post-step observable a ratio of two
affine functions of $u$, so the logarithmic sensitivity of
Section~\ref{sec:ground_fed_fraction} is $\mathrm{d}\ln\Robs/\mathrm{d}\ln u =
\ffrac{3}-\ffrac{4}$, the difference of the two ground-fed fractions.
Integrating between the frozen and the equilibrated reservoir gives, with no
approximation,
\begin{equation}
  \ln\frac{\Robs^{\rm PE}}{\Rcre}
  \;=\; \int_{\ln u^+}^{\ln u^-}\!\!\big(\ffrac{3}-\ffrac{4}\big)\,\mathrm{d}\ln u
  \;=\; \overline{S}\;\Delta\ln u ,
  \label{eq:exact_decomposition}
\end{equation}

\noindent where $\Delta\ln u = \ln(u^-/u^+)$ is the reservoir displacement and
\begin{equation}
  \overline{S} \;\equiv\;
  \frac{1}{\Delta\ln u}\int_{\ln u^+}^{\ln u^-}\!\!
  \big(\ffrac{3}-\ffrac{4}\big)\,\mathrm{d}\ln u
  \label{eq:Sbar_def}
\end{equation}

\noindent is the mean sensitivity over the interval the reservoir traverses.
Both factors are dimensionless. Since $\ffrac{3},\ffrac{4}\in[0,1]$ by
Section~\ref{sec:ground_fed_fraction}, and $\overline{S}$ is an average of their
difference, $|\overline{S}|$ is bounded by the peak sensitivity of
Eq.~\eqref{eq:peak_sensitivity} and so is strictly less than unity everywhere.
The plateau error follows:
\begin{equation}
  \boxed{\ \epsplat = \left|\,\mathrm{e}^{\overline{S}\,\Delta\ln u} - 1\,\right|\ }
  \label{eq:eps_plateau_exact}
\end{equation}

The consequence worth extracting is that a single scalar error has become a
product of two quantities that can be measured independently: how strongly the
observable responds to its reservoir, and how far the reservoir moved. Nothing
in Eq.~\eqref{eq:exact_decomposition} says which of them varies across the
operating map, and that is what the rest of this section determines.

\textbf{Use Eq.~\eqref{eq:eps_plateau_exact}, not its linearisation.} The
first-order form $\epsplat \approx |\overline{S}\,\Delta\ln u|$ is adequate for
a small step and wrong for a large one: at the four-interval step used below it
returns $59.4\,\%$ where the exact expression gives $81.1\,\%$, understating by
more than a quarter (\texttt{verify\_ridge\_mechanism.py}). The exponential
costs nothing.

One caution about what follows. Because Eq.~\eqref{eq:exact_decomposition} is an
identity, comparing its two sides tests nothing: a check that $\epsplat$ equals
$|\mathrm{e}^{\overline{S}\Delta\ln u}-1|$ can fail only through arithmetic
error. What the identity permits is an \emph{attribution} --- the two factors
are separately measurable, so one may ask which carries the structure the maps
show. That is a weaker kind of statement than a derivation, and it is the kind
made below.

\subsection{Across density: a sensitivity structure}
\label{sec:density_mechanism}

The plateau error has an interior maximum in density at
$\nel \approx 1.93\times10^{13}\,\si{\per\cubic\centi\metre}$, at every
temperature in the range this thesis defends.

Evaluating both factors of Eq.~\eqref{eq:exact_decomposition} at every grid
point settles which of them puts it there
(\texttt{verify\_ridge\_mechanism.py}, over $\Te = 2$--$9\,$eV and all eight
density columns). The column maximising $\epsplat$ is the column maximising
$|\overline{S}|$ in $32$ of $32$ temperature rows for a one-interval
temperature step and $31$ of $31$ for a four-interval step; the column
maximising $|\Delta\ln u|$ coincides with the error maximum in $0$ of $32$ and
$0$ of $31$.

Coincidence counts are weak evidence on their own, because a quantity that
barely varies still has its maximum somewhere. The amplitudes settle it. Across
the three decades of density $|\overline{S}|$ varies by a factor $7.11$ and
$|\Delta\ln u|$ by a factor $1.01$. \textbf{A quantity that changes by one per
cent cannot produce a ridge in which the error changes by a factor of seven},
wherever its own maximum falls. Table~\ref{tab:mechanism_profiles} makes the
point without statistics.

\begin{table}[htbp]
\centering
\caption{Density profiles of the plateau error and of the two factors of
Eq.~\eqref{eq:exact_decomposition}, each normalised to its own row maximum and
then taken as the median over the $32$ temperature rows with
$\Te = 2$--$9\,$eV, for a one-interval temperature step. From
\texttt{verify\_ridge\_mechanism.py}.}
\label{tab:mechanism_profiles}
\sisetup{table-format=1.3}
\begin{tabular}{lSSSSSSSS}
  \toprule
  $\nel$ (\si{\per\cubic\centi\metre}) & {$10^{12}$} & {$2.7{\cdot}10^{12}$} &
  {$7.2{\cdot}10^{12}$} & {$1.9{\cdot}10^{13}$} & {$5.2{\cdot}10^{13}$} &
  {$1.4{\cdot}10^{14}$} & {$3.7{\cdot}10^{14}$} & {$10^{15}$} \\
  \midrule
  $\epsplat$        & 0.437 & 0.683 & 0.937 & 1.000 & 0.807 & 0.512 & 0.280 & 0.138 \\
  $|\overline{S}|$  & 0.445 & 0.689 & 0.938 & 1.000 & 0.813 & 0.520 & 0.285 & 0.141 \\
  $|\Delta\ln u|$   & 0.997 & 0.995 & 0.993 & 0.991 & 0.990 & 0.992 & 0.996 & 1.000 \\
  \bottomrule
\end{tabular}
\end{table}

The first two rows agree to $2.2\,\%$ at every column; the third is flat to one
per cent. The density structure of the plateau error is the density structure of
the observable's sensitivity to its reservoir.

That sensitivity has an interior maximum for the reason
Section~\ref{sec:ground_fed_fraction} gives. Both $\ffrac{3}$ and $\ffrac{4}$
rise from zero to unity as the reservoir grows, but the two shells switch from
recombination-fed to ground-fed at different reservoir strengths, because the
ground-fed channel weakens with principal quantum number faster than the
recombination-fed one. Their difference therefore vanishes in both supply limits
and peaks between them. The density ridge is where the reservoir sits in that
mixed-supply region.

\subsection{Alternatives that would produce the same appearance}
\label{sec:ridge_alternatives}

An interior maximum on a coarse grid can be manufactured in several ways, and
the attribution above does not by itself exclude them.

\emph{Grid resolution.} The ridge occupies one density column and the grid is
spaced by a factor $2.683$. Its location is therefore known only to that factor,
and this work cannot say whether the true maximum lies at
$1.93\times10^{13}\,\si{\per\cubic\centi\metre}$ or anywhere within a factor
$2.7$ of it. What is established is that the maximum is interior --- not at
either end of the density range --- and that it does not move between
temperature rows. Refining the density grid requires rebuilding the rate matrix
and revalidating every number in Chapter~\ref{ch:theory} against the new one;
future work could determine the location more precisely.

\emph{Step-size convention.} A ridge whose position depended on the size of the
imposed step would be an artifact of that choice. It does not: the error maximum
falls on the same density column for a one-interval and a four-interval step in
$31$ of the $31$ temperature rows both can reach, and so does the maximum of
$|\overline{S}|$ (\texttt{verify\_ridge\_mechanism.py}). This tests steps of
fixed \emph{fractional} size; a step of fixed absolute size in electron-volts
spans a different number of grid intervals at each temperature and is a
different experiment, not tested here.

\emph{Truncation.} The model resolves $n\le8$ and bundles $n=9$--$15$. The
mechanism turns on the difference between the $n=3$ and $n=4$ supply fractions,
both fully resolved, but those fractions depend on the recombination flux
cascading from above, which the bundling approximates. This work does not test
the sensitivity of the ridge to $n_{\max}$; Chapter~\ref{ch:limits} lists it
among the outstanding checks.

\subsection{Across temperature: a displacement effect}
\label{sec:temperature_mechanism}

The temperature dependence has the opposite explanation, and this is where
treating the two axes alike would go wrong.

Along the ridge column the plateau error falls monotonically from $18.07\,\%$ at
$\Te = 2.02\,$eV to $4.55\,\%$ at $8.29\,$eV, a factor $3.72$. Decomposed
through Eq.~\eqref{eq:exact_decomposition}, $|\Delta\ln u|$ contributes a factor
$2.72$ of that fall and $|\overline{S}|$ only $1.29$; in logarithmic terms the
displacement carries $76\,\%$ of the decay and the sensitivity $19\,\%$. At the
four-interval step the split is $72\,\%$ and $13\,\%$
(\texttt{verify\_ridge\_mechanism.py}).

The observable therefore does not become much less sensitive as the plasma
heats. \textbf{The ground state simply moves less.} A step of fixed fractional
size displaces the reservoir progressively less at higher temperature, and the
error follows it down.

Why the displacement behaves this way is not established here. In particular
$|\Delta\ln u|$ is flat to one per cent across three decades of density while
varying by a factor $2.7$ across temperature, which implies that
$\partial\ln\ngs/\partial\ln\Te$ at collisional--radiative equilibrium is
controlled by temperature alone. \emph{[MECHANISM NOT ESTABLISHED: this is an
observation from \texttt{verify\_ridge\_mechanism.py}, not a derived result. It
has not been shown analytically and is a candidate for future work.]}

There is no interior maximum in temperature anywhere in the defended range. The
largest values, near $\Te\approx2\,$eV, are boundary values of the
optical-thickness restriction of Section~\ref{sec:optical_depth}, not a peak.

\subsection{Three explanations this work does not support}
\label{sec:withdrawn}

Three readings of the ridge are natural and do not survive the numbers. Each is
recorded because it is among the first a reader is likely to propose.

\textbf{It is not detachment.} The ridge sits at
$1.93\times10^{19}\,\si{\per\cubic\metre}$: a factor $5.2$ below the density
band that published edge modelling associates with a detached ITER divertor
\emph{[SOURCE REQUIRED: Guillemaut et al.\ 2011, Fus.\ Eng.\ Des.\ 86, 2954 ---
key not yet in the bibliography]}, and a factor $1.6$ below the modelled
upstream separatrix density \emph{[SOURCE REQUIRED: Stangeby et al.\ 2022,
Nucl.\ Fusion 62]}. It falls between the two, in a region neither source
characterises. The ridge is reproducible; its identification with a named
divertor operating regime is not supported by anything in this work.

\textbf{It is not a peak in temperature.} The error decreases monotonically with
temperature throughout the defended range, as
Section~\ref{sec:temperature_mechanism} shows.

\textbf{It is not simply where $\bdep^{\rm CRE}$ meets the sensitivity
maximum.} That reading accounts for the density structure but fails on the
temperature axis: at the four-interval step the ratio
$\bdep^{\rm CRE}/\bdep^{\rm peak}$ passes through unity at $\Te = 4.7\,$eV,
where the error is $33\,\%$ and falling monotonically
(\texttt{verify\_plateau\_gridmap.py}). The endpoint of the reservoir
trajectory is the wrong object: Eq.~\eqref{eq:Sbar_def} shows that what enters
is the mean sensitivity over the whole interval traversed.

\subsection{What this section establishes}

The plateau error factorises exactly into a sensitivity and a displacement,
Eq.~\eqref{eq:eps_plateau_exact}. Attributing the structure of the error maps to
one factor or the other gives different answers on the two axes: across density
the sensitivity carries the structure while the displacement is flat to one per
cent, and across temperature the displacement carries three quarters of the
decay. The ridge is a property of how the $n=3$ and $n=4$ shells divide their
supply between excitation from the ground state and recombination from the
continuum --- a property of the atomic system, needing no divertor geometry to
produce it --- while its amplitude at any density is set by how far a
temperature step moves the ground state.

That the local model generates a ridge without spatial structure does not imply
spatial structure is irrelevant to a real line-integrated measurement, which
Chapter~\ref{ch:limits} takes up.

The argument so far is entirely algebraic. Every quantity above is built from
the partial-equilibrium state $\Robs^{\rm PE}$, and that state was constructed
by freezing the ground state by hand rather than by integrating anything.
Whether the time-dependent system actually visits it is a separate question.


% ------------------------------------------------------------------
% Sections still to be written. The labels are declared here so that the
% cross-references in chapter5_C5D.tex resolve rather than printing "??".
% Replace each stub with the real section as it is drafted.
% ------------------------------------------------------------------
\section{Does the transient visit the plateau?}
\label{sec:bridge}

\emph{[TO BE WRITTEN --- C5-A writeup. Full 43-state integration; the
dynamic bridge; the closed form for the early-response slope $k$, which was
cut from Section~\ref{sec:mechanism} because it needs the transient
established first.]}

\section{How long the error persists}
\label{sec:persistence}

\emph{[TO BE WRITTEN. $\tau_{\rm eff}/\tauqss = 1.01$ and $1.07$.]}

\section{Optical thickness and the restricted range}
\label{sec:optical_depth}

\emph{[TO BE WRITTEN --- may instead live in Chapter~\ref{ch:theory};
Section~\ref{sec:mechanism} refers to it for the $\Te \ge 2$~eV restriction.]}
            % if they use \input
%
% Use whichever the neighbouring chapters use -- mixing them causes the page
% breaks and numbering to differ from the rest of the thesis.
% ============================================================

\chapter{What the Model Says}
\label{ch:results}

% ------------------------------------------------------------------
% PLACEHOLDER so the chapter is never empty while sections are drafted.
% Delete this block once sec:error_maps exists.
% ------------------------------------------------------------------
\section{The error maps}
\label{sec:error_maps}

\emph{[TO BE WRITTEN --- C5-F. Defines the operating grid, presents
$\epsstep$ and $\epsplat$ over it, and gives the headline numbers with their
scopes. Section~\ref{sec:mechanism} below is drafted and refers back to this
section for the maps it explains.]}

% ------------------------------------------------------------------
% C5-D, drafted 24 August 2026.
% ------------------------------------------------------------------
% ============================================================
% Chapter 5, Section: the mechanism.  C5-D, draft 2.
%
% Written under the thesis-writing skill, 24 August 2026.
% Source: CH5_EVIDENCE.md. L_grid SHA 2d92b58e..., S_grid SHA 7822f536...
% Draft 1 retained as chapter5_C5D_draft1_20260824.tex (failed six gates).
%
% MACROS beyond those chapter3.tex exercises:  \epsstep  \epsplat
%   (recorded in notation_and_definitions.md Sec. 7, never yet compiled --
%    CONFIRM they exist in thesis_main.tex before building.)
% ============================================================

%-----------------------------------------------------------------------------
\section{Why the error is largest where it is}
\label{sec:mechanism}
%-----------------------------------------------------------------------------

The maps of Section~\ref{sec:error_maps} show where the diagnostic error is
large. They do not say why, and a map without a mechanism is a lookup table
that cannot be extrapolated past the grid that produced it. This section asks
which physical quantity puts the ridge where it is.

The answer has two halves and they disagree. The structure along the density
axis and the structure along the temperature axis have different causes, and
reading them as one effect --- which the appearance of a single ridged surface
invites --- gets both wrong.

\subsection{Two error measures}
\label{sec:eps_definitions}

Section~\ref{sec:two_references} fixed the two reference states against which an
observable can be judged. Two further quantities are needed here, and both are
properties of a temperature \emph{step} rather than of a single operator, which
is why neither appears in Chapter~\ref{ch:theory}.

Let the plasma sit at the collisional--radiative equilibrium of
$\Lmat(\Te,\nel)$, and let the electron temperature step to $\Te+\Delta\Te$ at
$t=0$ with $\nel$ held fixed. Write $\Lmat^-$ and $\Lmat^+$ for the operators
before and after, in the notation of Section~\ref{sec:two_clocks}. The
\emph{step error}
\begin{equation}
  \epsstep = \left|\frac{\Rcre(\Te,\nel)}
                        {\Rcre(\Te+\Delta\Te,\nel)} - 1\right|
  \label{eq:eps_step_def}
\end{equation}

\noindent is the dimensionless distance between the two tabulated answers: what
the reading would change by if the plasma re-equilibrated instantly. The
\emph{plateau error}
\begin{equation}
  \epsplat = \left|\frac{\Robs^{\rm PE}}
                        {\Rcre(\Te+\Delta\Te,\nel)} - 1\right|,
  \qquad
  \Robs^{\rm PE} \equiv \Rqss\!\left(\Te+\Delta\Te,\nel;\,\bdep^-\right),
  \label{eq:eps_plateau_def}
\end{equation}

\noindent is the distance between the tabulated answer and the state the plasma
occupies once the excited manifold has settled but before the ground state has
moved. The superscript PE denotes \emph{partial equilibrium}: the post-step
excited response of Eq.~\eqref{eq:qss_ratio}, evaluated on the \emph{pre-step}
reservoir.

The consequence of keeping these separate is that they are not close. Across the
operating grid $\epsplat$ exceeds $\epsstep$ at every point tested, typically by
an order of magnitude (Section~\ref{sec:error_maps}). A diagnostician who
pictures the reading moving smoothly between the two tabulated values has not
merely misjudged the size of the excursion but has the wrong shape of transient.
$\epsplat$ is the quantity this chapter reports, and Section~\ref{sec:bridge}
establishes that the time-dependent system does pass through the state it
measures.

\subsection{An exact decomposition of the finite step}
\label{sec:exact_decomposition}

Chapter~\ref{ch:theory} worked in the ground-state departure coefficient
$\bdep = \ngs/(\Zsb\nel\nion)$, because that is the variable in which the
observable's response has the closed form of Eq.~\eqref{eq:R_of_b}. Here it is
better to use the reservoir itself,
\begin{equation}
  u \;\equiv\; \frac{\ngs}{\nion} \;=\; \bdep\,\Zsb(\Te)\,\nel ,
  \label{eq:u_def}
\end{equation}

\noindent a dimensionless ground-state population per unit ion density. At
fixed $(\Te,\nel)$ the two are proportional, so $\mathrm{d}\ln u =
\mathrm{d}\ln\bdep$ and the sensitivity of
Section~\ref{sec:ground_fed_fraction} carries over unchanged. The reason for
the change of variable is that a temperature step moves $\Zsb$: comparing a
pre-step $\bdep$ with a post-step one would mix a change in the reservoir with
a change in the Saha--Boltzmann normalisation, and $u$ has no normalisation to
move. That confusion is not hypothetical --- an earlier version of the analysis
reported in this section made exactly it, and the resulting contamination in
$\Delta\ln\bdep$ was comparable to the quantity being measured.

\textbf{Assumption.} Equation~\eqref{eq:u_def} treats $\nion$ as a fixed
reservoir, unchanged by the step, as throughout this thesis
(Section~\ref{sec:rate_equation}). This removes the freedom for the ionisation
balance to respond, and it is what makes $u$ numerically equal to the
ground-state population in the solved system. Chapter~\ref{ch:limits} returns
to what a mobile ion population would change.

With that, Eq.~\eqref{eq:R_of_b} makes the post-step observable a ratio of two
affine functions of $u$, so the logarithmic sensitivity of
Section~\ref{sec:ground_fed_fraction} is $\mathrm{d}\ln\Robs/\mathrm{d}\ln u =
\ffrac{3}-\ffrac{4}$, the difference of the two ground-fed fractions.
Integrating between the frozen and the equilibrated reservoir gives, with no
approximation,
\begin{equation}
  \ln\frac{\Robs^{\rm PE}}{\Rcre}
  \;=\; \int_{\ln u^+}^{\ln u^-}\!\!\big(\ffrac{3}-\ffrac{4}\big)\,\mathrm{d}\ln u
  \;=\; \overline{S}\;\Delta\ln u ,
  \label{eq:exact_decomposition}
\end{equation}

\noindent where $\Delta\ln u = \ln(u^-/u^+)$ is the reservoir displacement and
\begin{equation}
  \overline{S} \;\equiv\;
  \frac{1}{\Delta\ln u}\int_{\ln u^+}^{\ln u^-}\!\!
  \big(\ffrac{3}-\ffrac{4}\big)\,\mathrm{d}\ln u
  \label{eq:Sbar_def}
\end{equation}

\noindent is the mean sensitivity over the interval the reservoir traverses.
Both factors are dimensionless. Since $\ffrac{3},\ffrac{4}\in[0,1]$ by
Section~\ref{sec:ground_fed_fraction}, and $\overline{S}$ is an average of their
difference, $|\overline{S}|$ is bounded by the peak sensitivity of
Eq.~\eqref{eq:peak_sensitivity} and so is strictly less than unity everywhere.
The plateau error follows:
\begin{equation}
  \boxed{\ \epsplat = \left|\,\mathrm{e}^{\overline{S}\,\Delta\ln u} - 1\,\right|\ }
  \label{eq:eps_plateau_exact}
\end{equation}

The consequence worth extracting is that a single scalar error has become a
product of two quantities that can be measured independently: how strongly the
observable responds to its reservoir, and how far the reservoir moved. Nothing
in Eq.~\eqref{eq:exact_decomposition} says which of them varies across the
operating map, and that is what the rest of this section determines.

\textbf{Use Eq.~\eqref{eq:eps_plateau_exact}, not its linearisation.} The
first-order form $\epsplat \approx |\overline{S}\,\Delta\ln u|$ is adequate for
a small step and wrong for a large one: at the four-interval step used below it
returns $59.4\,\%$ where the exact expression gives $81.1\,\%$, understating by
more than a quarter (\texttt{verify\_ridge\_mechanism.py}). The exponential
costs nothing.

One caution about what follows. Because Eq.~\eqref{eq:exact_decomposition} is an
identity, comparing its two sides tests nothing: a check that $\epsplat$ equals
$|\mathrm{e}^{\overline{S}\Delta\ln u}-1|$ can fail only through arithmetic
error. What the identity permits is an \emph{attribution} --- the two factors
are separately measurable, so one may ask which carries the structure the maps
show. That is a weaker kind of statement than a derivation, and it is the kind
made below.

\subsection{Across density: a sensitivity structure}
\label{sec:density_mechanism}

The plateau error has an interior maximum in density at
$\nel \approx 1.93\times10^{13}\,\si{\per\cubic\centi\metre}$, at every
temperature in the range this thesis defends.

Evaluating both factors of Eq.~\eqref{eq:exact_decomposition} at every grid
point settles which of them puts it there
(\texttt{verify\_ridge\_mechanism.py}, over $\Te = 2$--$9\,$eV and all eight
density columns). The column maximising $\epsplat$ is the column maximising
$|\overline{S}|$ in $32$ of $32$ temperature rows for a one-interval
temperature step and $31$ of $31$ for a four-interval step; the column
maximising $|\Delta\ln u|$ coincides with the error maximum in $0$ of $32$ and
$0$ of $31$.

Coincidence counts are weak evidence on their own, because a quantity that
barely varies still has its maximum somewhere. The amplitudes settle it. Across
the three decades of density $|\overline{S}|$ varies by a factor $7.11$ and
$|\Delta\ln u|$ by a factor $1.01$. \textbf{A quantity that changes by one per
cent cannot produce a ridge in which the error changes by a factor of seven},
wherever its own maximum falls. Table~\ref{tab:mechanism_profiles} makes the
point without statistics.

\begin{table}[htbp]
\centering
\caption{Density profiles of the plateau error and of the two factors of
Eq.~\eqref{eq:exact_decomposition}, each normalised to its own row maximum and
then taken as the median over the $32$ temperature rows with
$\Te = 2$--$9\,$eV, for a one-interval temperature step. From
\texttt{verify\_ridge\_mechanism.py}.}
\label{tab:mechanism_profiles}
\sisetup{table-format=1.3}
\begin{tabular}{lSSSSSSSS}
  \toprule
  $\nel$ (\si{\per\cubic\centi\metre}) & {$10^{12}$} & {$2.7{\cdot}10^{12}$} &
  {$7.2{\cdot}10^{12}$} & {$1.9{\cdot}10^{13}$} & {$5.2{\cdot}10^{13}$} &
  {$1.4{\cdot}10^{14}$} & {$3.7{\cdot}10^{14}$} & {$10^{15}$} \\
  \midrule
  $\epsplat$        & 0.437 & 0.683 & 0.937 & 1.000 & 0.807 & 0.512 & 0.280 & 0.138 \\
  $|\overline{S}|$  & 0.445 & 0.689 & 0.938 & 1.000 & 0.813 & 0.520 & 0.285 & 0.141 \\
  $|\Delta\ln u|$   & 0.997 & 0.995 & 0.993 & 0.991 & 0.990 & 0.992 & 0.996 & 1.000 \\
  \bottomrule
\end{tabular}
\end{table}

The first two rows agree to $2.2\,\%$ at every column; the third is flat to one
per cent. The density structure of the plateau error is the density structure of
the observable's sensitivity to its reservoir.

That sensitivity has an interior maximum for the reason
Section~\ref{sec:ground_fed_fraction} gives. Both $\ffrac{3}$ and $\ffrac{4}$
rise from zero to unity as the reservoir grows, but the two shells switch from
recombination-fed to ground-fed at different reservoir strengths, because the
ground-fed channel weakens with principal quantum number faster than the
recombination-fed one. Their difference therefore vanishes in both supply limits
and peaks between them. The density ridge is where the reservoir sits in that
mixed-supply region.

\subsection{Alternatives that would produce the same appearance}
\label{sec:ridge_alternatives}

An interior maximum on a coarse grid can be manufactured in several ways, and
the attribution above does not by itself exclude them.

\emph{Grid resolution.} The ridge occupies one density column and the grid is
spaced by a factor $2.683$. Its location is therefore known only to that factor,
and this work cannot say whether the true maximum lies at
$1.93\times10^{13}\,\si{\per\cubic\centi\metre}$ or anywhere within a factor
$2.7$ of it. What is established is that the maximum is interior --- not at
either end of the density range --- and that it does not move between
temperature rows. Refining the density grid requires rebuilding the rate matrix
and revalidating every number in Chapter~\ref{ch:theory} against the new one;
future work could determine the location more precisely.

\emph{Step-size convention.} A ridge whose position depended on the size of the
imposed step would be an artifact of that choice. It does not: the error maximum
falls on the same density column for a one-interval and a four-interval step in
$31$ of the $31$ temperature rows both can reach, and so does the maximum of
$|\overline{S}|$ (\texttt{verify\_ridge\_mechanism.py}). This tests steps of
fixed \emph{fractional} size; a step of fixed absolute size in electron-volts
spans a different number of grid intervals at each temperature and is a
different experiment, not tested here.

\emph{Truncation.} The model resolves $n\le8$ and bundles $n=9$--$15$. The
mechanism turns on the difference between the $n=3$ and $n=4$ supply fractions,
both fully resolved, but those fractions depend on the recombination flux
cascading from above, which the bundling approximates. This work does not test
the sensitivity of the ridge to $n_{\max}$; Chapter~\ref{ch:limits} lists it
among the outstanding checks.

\subsection{Across temperature: a displacement effect}
\label{sec:temperature_mechanism}

The temperature dependence has the opposite explanation, and this is where
treating the two axes alike would go wrong.

Along the ridge column the plateau error falls monotonically from $18.07\,\%$ at
$\Te = 2.02\,$eV to $4.55\,\%$ at $8.29\,$eV, a factor $3.72$. Decomposed
through Eq.~\eqref{eq:exact_decomposition}, $|\Delta\ln u|$ contributes a factor
$2.72$ of that fall and $|\overline{S}|$ only $1.29$; in logarithmic terms the
displacement carries $76\,\%$ of the decay and the sensitivity $19\,\%$. At the
four-interval step the split is $72\,\%$ and $13\,\%$
(\texttt{verify\_ridge\_mechanism.py}).

The observable therefore does not become much less sensitive as the plasma
heats. \textbf{The ground state simply moves less.} A step of fixed fractional
size displaces the reservoir progressively less at higher temperature, and the
error follows it down.

Why the displacement behaves this way is not established here. In particular
$|\Delta\ln u|$ is flat to one per cent across three decades of density while
varying by a factor $2.7$ across temperature, which implies that
$\partial\ln\ngs/\partial\ln\Te$ at collisional--radiative equilibrium is
controlled by temperature alone. \emph{[MECHANISM NOT ESTABLISHED: this is an
observation from \texttt{verify\_ridge\_mechanism.py}, not a derived result. It
has not been shown analytically and is a candidate for future work.]}

There is no interior maximum in temperature anywhere in the defended range. The
largest values, near $\Te\approx2\,$eV, are boundary values of the
optical-thickness restriction of Section~\ref{sec:optical_depth}, not a peak.

\subsection{Three explanations this work does not support}
\label{sec:withdrawn}

Three readings of the ridge are natural and do not survive the numbers. Each is
recorded because it is among the first a reader is likely to propose.

\textbf{It is not detachment.} The ridge sits at
$1.93\times10^{19}\,\si{\per\cubic\metre}$: a factor $5.2$ below the density
band that published edge modelling associates with a detached ITER divertor
\emph{[SOURCE REQUIRED: Guillemaut et al.\ 2011, Fus.\ Eng.\ Des.\ 86, 2954 ---
key not yet in the bibliography]}, and a factor $1.6$ below the modelled
upstream separatrix density \emph{[SOURCE REQUIRED: Stangeby et al.\ 2022,
Nucl.\ Fusion 62]}. It falls between the two, in a region neither source
characterises. The ridge is reproducible; its identification with a named
divertor operating regime is not supported by anything in this work.

\textbf{It is not a peak in temperature.} The error decreases monotonically with
temperature throughout the defended range, as
Section~\ref{sec:temperature_mechanism} shows.

\textbf{It is not simply where $\bdep^{\rm CRE}$ meets the sensitivity
maximum.} That reading accounts for the density structure but fails on the
temperature axis: at the four-interval step the ratio
$\bdep^{\rm CRE}/\bdep^{\rm peak}$ passes through unity at $\Te = 4.7\,$eV,
where the error is $33\,\%$ and falling monotonically
(\texttt{verify\_plateau\_gridmap.py}). The endpoint of the reservoir
trajectory is the wrong object: Eq.~\eqref{eq:Sbar_def} shows that what enters
is the mean sensitivity over the whole interval traversed.

\subsection{What this section establishes}

The plateau error factorises exactly into a sensitivity and a displacement,
Eq.~\eqref{eq:eps_plateau_exact}. Attributing the structure of the error maps to
one factor or the other gives different answers on the two axes: across density
the sensitivity carries the structure while the displacement is flat to one per
cent, and across temperature the displacement carries three quarters of the
decay. The ridge is a property of how the $n=3$ and $n=4$ shells divide their
supply between excitation from the ground state and recombination from the
continuum --- a property of the atomic system, needing no divertor geometry to
produce it --- while its amplitude at any density is set by how far a
temperature step moves the ground state.

That the local model generates a ridge without spatial structure does not imply
spatial structure is irrelevant to a real line-integrated measurement, which
Chapter~\ref{ch:limits} takes up.

The argument so far is entirely algebraic. Every quantity above is built from
the partial-equilibrium state $\Robs^{\rm PE}$, and that state was constructed
by freezing the ground state by hand rather than by integrating anything.
Whether the time-dependent system actually visits it is a separate question.


% ------------------------------------------------------------------
% Sections still to be written. The labels are declared here so that the
% cross-references in chapter5_C5D.tex resolve rather than printing "??".
% Replace each stub with the real section as it is drafted.
% ------------------------------------------------------------------
\section{Does the transient visit the plateau?}
\label{sec:bridge}

\emph{[TO BE WRITTEN --- C5-A writeup. Full 43-state integration; the
dynamic bridge; the closed form for the early-response slope $k$, which was
cut from Section~\ref{sec:mechanism} because it needs the transient
established first.]}

\section{How long the error persists}
\label{sec:persistence}

\emph{[TO BE WRITTEN. $\tau_{\rm eff}/\tauqss = 1.01$ and $1.07$.]}

\section{Optical thickness and the restricted range}
\label{sec:optical_depth}

\emph{[TO BE WRITTEN --- may instead live in Chapter~\ref{ch:theory};
Section~\ref{sec:mechanism} refers to it for the $\Te \ge 2$~eV restriction.]}
